% The next line makes it so it compiles the main file if I hit the compile button
% when editing this file in TeXworks.
% !TEX root = ../AIDMA.tex
%%%%%%%%%%%%%%%%%%%%%%%%%%%%%%%%%%%%

% The newline causes a warning about \\ being invalid in a PDF string, but
% I get an overfull hbox without it.
\chapter{Recursion, Recurrences, and\\ 
Mathematical Induction}\label{ch:recurrences}
In this chapter we will explore 
a proof technique, an algorithmic technique, and a mathematical technique.  
Each topic is in some ways very different than the others, 
yet they have a whole lot in common.  They are also often used in conjunction.

You have already seen {\em recurrence relations}.  Recall that a recurrence relation is a way of
defining a sequence of numbers with a formula that is based on previous numbers in the sequence.
You are probably also familiar with {\em recursion},
an algorithmic technique in which
an algorithm calls itself (such an algorithm is called {\em recursive}), typically with ``smaller'' input.  
Finally, the {\em principle of mathematical induction} is a slick 
proof technique that works so well that sometimes it feels like 
you are cheating.

We will see that induction can be used to prove formulas, 
prove that algorithms---especially recursive ones---are correct, and help solve recurrence
relations.  Among other things, recurrence relations can be used to analyze recursive algorithm.  
Recursive algorithms can be used to compute the values defined by recurrence
relations and to solve problems that can be broken into smaller versions of themselves.

As we will see, each of these has one or more {\em base cases} that can be proved/computed/determined
directly and a {\em recursive} or {\em inductive} step that relies on previous steps.
With each, the inductive/recursive steps must eventually lead to a base case.

Because induction can be used to prove things about the other two, we will begin there.

\section{Mathematical Induction}\label{sec:induction}
\index{induction}
\index{mathematical induction}
Let's begin our study of mathematical induction (often just called induction) with an example that should look familiar.  
It is actually Theorem~\ref{thm:powerSetSize} that we proved 
in an earlier chapter.
Following that, we will explain how/why induction works and give
plenty of other examples.


\begin{exa}\label{induction1}
Let $A$ be a set with $n$ elements.  Prove that $|P(A)|=2^n$. 
\begin{pf}
We use induction and the idea from the solution to Exercise~\ref{exa:subsets4element}. 
Clearly if $|A|= 1$, $A$
has $2^1 = 2$ subsets: $\varnothing$ and $A$ itself.\\
Assume every
set with $n - 1$ elements has $2^{n-1}$ subsets. Let $A$ be a set
with $n$ elements.  Choose some $x\in A$.  Every subset of $A$
either contains $x$ or it doesn't.  
Those that do not contain $x$  are subsets of $A\setminus \{x\}$.  Since
$A\setminus \{x\}$ has $n - 1$ elements, the induction hypothesis implies that it has
$2^{n-1}$ subsets.
Every subset that does contain $x$ corresponds to one of the subsets of $A\setminus \{x\}$ with the element $x$ added.  That is, for each subset $S \subseteq A\setminus \{x\}$,  $S\cup \{x\}$ is a subset of $A$ containing $x$.
Clearly there are $2^{n-1}$ such new subsets.
Since this accounts for all subsets of $A$, $A$ has $2^{n-1} + 2^{n-1} = 2^n$ subsets.  
\end{pf}
\end{exa}

Now we will go into detail about how and why induction works.
You should come back and reread Example~\ref{induction1} 
after reading section~\ref{sec:indBasic}.

\subsection{The Basics}\label{sec:indBasic}
The {\em principle of mathematical induction} 
\index{mathematical induction}
\index{induction}
\index{proof!induction}
({\em PMI}, or simply {\em induction}) is usually used to 
prove statements of the form 
$$\mbox{for all } n\geq a, P(n) \mbox{ is true},$$  
where $a$ is an integer,
and $P(n)$ is a propositional function with domain $\{a, a+1, a+2,\ldots\}$. 
Usually $a$ is $0$ or $1$,
so the domain is usually $\naturals$ (the natural numbers) or $\BBZ^+$
(the positive integers).


Induction is based on the following
fairly intuitive observation (which we will formalize next). Suppose that we are to perform a
task that involves a certain number of steps. Suppose that these
steps must be followed in strict numerical order. Finally, suppose
that we know how to perform the $n$-th task provided we have
accomplished the $(n - 1)$-th task. Thus if we are ever able to
start the job (that is, if we have a base case), then we should be
able to finish it (because starting with the base case we go to
the next case, and then to the case following that, etc.).

\begin{exer}
Based on the description so far, which of the following
statements {\em might} we be able to prove with mathematical induction (indicate with `Y' or `N')?
Briefly justify.
\begin{lenum}
\itemtf
The square of any integer is positive.

\vspace*{.125in}
\itemtf
Every positive integer can be written as the sum of two other
positive integers.

\vspace*{.125in}
\itemtf
Every integer greater than $1$ can be written as the product of
prime numbers.

\vspace*{.125in}
\itemtf
If $n\geq 1$, 
$\dsum _{k=1}^n k^2 = \dfrac{n(n+1)(2n+1)}{6}$

\vspace*{.125in}
\itemtf
Every real number is the square of another real number.

\end{lenum}

\vspace*{.25in} % After the enum because it doesn't work inside it for some reason.
\begin{answer}
(a) No.  The domain is $\BBZ$, which does not have a `starting point'.
(b) Yes.  The domain is $\BBZ^+$.
(c) Yes.  The domain is $\{2,3,4,\ldots\}$.
(d) Yes.  The domain is $\BBZ^+$.
(e) No.  
The domain is $\reals$ which is not a subset of $\BBZ$.  
Thus, not only is there no `starting point,'
there is no clear ordering of the real numbers from one to the next.
\end{answer}
\end{exer}

The following example illustrates the idea behind induction.
It uses {\em modus ponens}\index{modus ponens}.
Recall that modus ponens states that 
if $p$ is true and $p\rightarrow q$ is true, then $q$ is true.  In English, 
``If $p$ is true, and whenever $p$ is true $q$ is true, then $q$ is true.''\footnote{We can also write
this as the tautology $[p\wedge(p\rightarrow q)]\rightarrow q$.}

\begin{exa}\label{exa:ind1}
Assume that we know that $P(1)$ is true and that 
whenever $k\geq 1$, $P(k)\rightarrow P(k+1)$ is true.  What can we conclude?
\begin{sol}
Let's start from the ground up.
We know that $P(1)$ is true.  We also know that $P(k)\rightarrow P(k+1)$ is true for 
any integer $k\geq 1$.  
For instance, since $4\geq 1$, we know that $P(4)\rightarrow P(5)$ is true.  
It should be noted that we don't (yet) know anything about the truth values of 
$P(4)$ and $P(5)$.
\begin{itemize}
\item 
We know $P(1)$ is true and $1\geq 1$, $P(1)\rightarrow P(2)$ is true, therefore $P(2)$ is true.
\item
Since $P(2)$ is true and $2\geq 1$, $P(2)\rightarrow P(3)$ is true, therefore $P(3)$ is true.
\item
Since $P(3)$ is true and $3\geq 1$, $P(3)\rightarrow P(4)$ is true, therefore $P(4)$ is true.
\item
Since $P(4)$ is true and $4\geq 1$, $P(4)\rightarrow P(5)$ is true, therefore $P(5)$ is true.
\item
Since $P(5)$ is true and $5\geq 1$, $P(5)\rightarrow P(6)$ is true, therefore $P(6)$ is true.
\end{itemize}
It seems pretty clear that this pattern continues for all values of $k>6$ as well, 
so $P(k)$ is true for all $k\geq 1$.
\end{sol}
\end{exa}

\begin{ques}
Example~\ref{exa:ind1} had several statements like the following:
\begin{quote}
``Since $P(4)$ is true and $4\geq 1$, $P(4)\rightarrow P(5)$ is true, therefore $P(5)$ is true.''
\end{quote}
What is the justification for the conclusion that $P(5)$ is true?

\noindent
\answerline
\begin{answer}
Modus ponens.
\end{answer}
\end{ques}

Example~\ref{exa:ind1} did not give a formal
{\em proof} of the conclusion. 
The idea is to get you thinking
about how mathematical induction works, 
not to provide a formal proof that it does (yet).
Hopefully this example will help prime your brain for the proof that 
mathematical induction is a valid proof technique that we will 
give shortly.
%Once you wrap your head around mathematical induction (it takes some people
%longer than others), you will believe it works regardless of whether
%or not you remember the details of the proof.

Before moving on, we should make sure you understand
what has already been said.

\begin{ques}
If you know that $P(5)$ is true, and you also know that
$P(k)\rightarrow P(k+1)$ whenever $k\geq 1$,
what can you conclude?

\answerline
\begin{answer}
You can immediately conclude that $P(6)$ is true using
modus ponens.  If that was your answer, good.
But you can keep going.  Since $P(6)$ is true,
you can conclude that $P(7)$ is true (also by modus ponens).
But then you can conclude that $P(8)$ is true.  
And so on.  The most complete answer you can give is that
$P(n)$ is true for all $n\geq 5$.
You {\em cannot} conclude that $P(n)$ is true for
all $n\geq 1$ because we don't know anything about
the truth values of $P(1)$, $P(2)$, $P(3)$, and $P(4)$.
\end{answer}
\end{ques}

\begin{ques}
If you know that $P(17)$ is true and you also know that
$P(k)\rightarrow P(k+1)$ whenever $k\geq 1$, what can you conclude
about $P(10)$?

\answerline
\begin{answer}
Nothing. We can conclude that $P(n)$ is true for any $n\geq 17$,
but there is not enough information to say anything about values of $n$
less than 17.
\end{answer}
\end{ques}

Now it is time to get more formal with our discussion.
Mathematical induction is based on the fact that
if $P(a)$ is true for some $a\geq 0$ (the {\em base case}), 
\index{base case (induction)}
and for any $k\geq a$, if $P(k)$ is true, then $P(k+1)$ is true
(the {\em inductive case}),
then $P(n)$ is true for all $n\geq a$. 
In other words, the principle of mathematical induction is based on the
fact that 
$$[P(a)\wedge\forall k(P(k)\rightarrow P(k+1))]\rightarrow (\forall n P(n)),$$
where the universe is $\{a,a+1,a+2,\ldots\}$, is true.

\begin{exer}
Restate $[P(a)\wedge\forall k(P(k)\rightarrow P(k+1))]\rightarrow (\forall n P(n))$ 
(where the universe is $\{a,a+1,a+2,\ldots\}$) in English.

\noindent
\answerlinethree
\begin{answer}
There are various ways to say this, including what was said in the paragraph above.
Here is another way to say it:
\begin{quote}
If $P(a)$ is true, and for any value of $k\geq a$,
$P(k)$ true implies that $P(k+1)$ is true, then $P(n)$ is true for all $n\geq a$.
\end{quote}
\end{answer}
\end{exer}

%We won't prove that this is a tautology, but 
%hopefully 
The proof that $[P(a)\wedge\forall k(P(k)\rightarrow P(k+1))]\rightarrow (\forall n P(n))$ is true is based on something called the 
{\em well-ordering principle} which states that
every nonempty subset of the natural numbers has a least element.
Read the following proof very carefully, making sure you 
understand the justification of every step. If you are not sure about
any of the steps, it is important that you get them clarified!

\begin{thm}
Assume we are working over the universe $\{a, a+1, a+2, \ldots\}$.
The statement $[P(a)\wedge\forall k(P(k)\rightarrow P(k+1))]\rightarrow (\forall n P(n))$ is true.
\begin{pf}
If the statement is false, then it must be that
$P(a)\wedge\forall k(P(k)\rightarrow P(k+1))$
is true but that $\forall n P(n)$ is false.
Let $S=\{s\in \{a, a+1, a+2, \ldots\} | \neg P(s)\}$.
That is, $S$ is the set of integers for which $P(n)$ is false.
% be the set of integers such that $s\in S$ iff $P(s)$ is false. 
Since $\forall n P(n)$ is false, $S$ is nonempty.
Clearly $S$ is a subset of the natural numbers, so the well-ordering principle applies. Therefore there is some least element $b\in S$.
Since $b\in S$, $P(b)$ is false, and since it is the
least such element, $b-1\not\in S$, so $P(b-1)$ is true. 
But we know that $\forall k(P(k)\rightarrow P(k+1))$ is true, so
$P(b-1)\rightarrow P(b)$.
By modus ponens, $P(b)$ is true, a contradiction.
Therefore the statement is true.
\end{pf}
\end{thm}

%Because this is a bit technical, 
%we will omit a discussion of it along with a proof of the statement. 
%However, hopefully
%Example~\ref{exa:ind1} helped convince
%you that it is indeed a true.  
It is definitely worth your time to convince yourself that
mathematical induction is a valid technique.
If you aren't convinced, reread the proof, think about it some more, 
and/or ask someone to help you understand it.

\begin{ques}
Are you convinced that $[P(a)\wedge\forall k(P(k)\rightarrow P(k+1))]\rightarrow (\forall n P(n))$ is true?

\noindent
\answerline
\begin{answer}
If you answered {\bf yes} and you aren't lying, great!  If you
answered {\em no} or you answered {\bf yes} but you lied, 
it is important that you think about it some more and/or
get some help.  If you want to succeed at writing induction
proofs, understanding this is an important step!
\end{answer}
\end{ques}


We call $P(a)$ the {\em base case}.\index{base case}
Sometimes we actually need to prove several
base cases (we will see why later).  For instance,
we might need to prove $P(a)$, $P(a+1)$, and $P(a+2)$ 
are all true.

The {\em inductive step}\index{inductive step}
involves proving that  $\forall k(P(k)\rightarrow P(k+1))$
is true.  To prove it, we show that if $P(k)$ is true for any $k$ 
which is {\em at least as large as the base case(s)}, 
then $P(k+1)$ is true. 
The assumption that $P(k)$ is true is called the {\em inductive hypothesis}.\index{inductive hypothesis}

Based on our discussion so far, here is the procedure for writing induction proofs.

\newpage % A lot of empty space before here, but this should be on one page.

\begin{proc}\label{proc:induct1}
To use induction to prove that $\forall n P(n)$ is true on domain
$\{a,a+1,\ldots\}$:
\begin{enumerate}
\item
{\bf Base Case:} Show that $P(a)$ is true (and possible one or more additional base cases).
\item 
Show that $\forall k(P(k)\rightarrow P(k+1))$ is true.
To show this:
\begin{lenum}
\item
{\bf Inductive Hypothesis:} 
Let $k\geq a$ be an integer and assume that $P(k)$ is true.
\item
{\bf Inductive Step:}
Prove that $P(k+1)$ is true, typically using the fact that
$P(k)$ is true.
\end{lenum}
Assuming we used no special facts about $k$ other than 
$k\geq a$, this means we have shown that 
$\forall k(P(k)\rightarrow P(k+1))$ 
(again, where it is understood that the domain is $\{a,a+1,\ldots\})$.
\item
{\bf Summary:} Conclude that $\forall n P(n)$ is true, usually by saying
something like ``Since $P(a)$ and $P(k)\rightarrow P(k+1)$ for all $k\geq a$, 
$\forall n P(n)$ is true by induction.''
\end{enumerate}
\end{proc}

As you will quickly learn, the {\em base case} is generally pretty easy, 
as is writing down the {\em inductive hypothesis}. The {\em summary} is 
even easier, since it almost always says the same thing.
The {\em inductive step} is the longest and most complicated
step.  In fact, in mathematics and theoretical computer science journals, induction
proofs often only include the inductive step since anyone reading papers in such journals
can generally fill in the details of the other three parts.
But keep in mind that you are not (yet) writing papers for such journals, so you
{\em cannot} omit these steps!

Let's see another example. 



\begin{exa}\label{exa:sumodd}
Prove that the sum of the first $n$ odd integers is $n^2$. 
That is, show that 
$\dsum_{i=1}^{n}(2 i - 1) = n^2$ for all $n\geq 1$.
\begin{pf}
Let $P(n)$  be the statement  ``$\dsum_{i=1}^{n}(2 i - 1) = n^2$''.
We need to show that $P(n)$ is true for all $n\geq 1$.

\noindent
{\bf Base Case:} 
Since $\dsum_{i=1}^{1}(2 i - 1) = 2\cdot 1 - 1 = 1=1^2,$
$P(1)$ is true. 

\noindent
{\bf Inductive Hypothesis:}
Let $k\geq 1$ and assume that $P(k)$ is true.  That is,
assume that $\dsum_{i=1}^{k}(2 i - 1) = k^2$ when $k\geq 1$. \hfill 

\noindent
{\bf Inductive Step:} 
Then \hfill
\begin{eqnarray*}
\displaystyle \sum_{i=1}^{k+1} (2 i - 1)&=&
\displaystyle \sum_{i=1}^{k}( 2 i - 1)+ (2(k+1) - 1) \, \, \, 
\mbox{ (take }k+1\mbox{ term from sum)}\\
&=&k^2+(2k+2-1)\, \, \, 
\mbox{ (by the {\em inductive hypothesis})}\\
&=&k^2+2k+1\\
&=&(k+1)^2
\end{eqnarray*}
Thus $P(k+1)$ is true.

\noindent
{\bf Summary:}
Since we proved that $P(1)$ is true, and that $P(k)\rightarrow P(k+1)$ whenever $k\geq 1$,
$P(n)$ is true for all $n\geq 1$ by the principle of mathematical induction.
\end{pf}
\end{exa}

The previous proof had the four components we discussed.
We proved the {\em base case}.
We then assumed it was true for $k$.  That is, we made the {\em inductive hypothesis}.
Next we proved that it was true for $k+1$ based on the assumption that it is true for $k$.  
That is, we did the {\em inductive step}.
Finally, we appealed to the principle of mathematical induction in the {\em summary}.

\begin{rem}
Recall the following statement from 
 Example~\ref{exa:sumodd}:
\begin{quote}
Let $P(n)$  be the statement  ``$\dsum_{i=1}^{n}(2 i - 1) = n^2$''.
\end{quote} 
Did you notice the quotes?  
It is important that you include these.  This is particularly important if
you use notation such as  $P(n)=$``$\dsum_{i=1}^{n}(2 i - 1) = n^2$''.  Without the quotes, 
this becomes $P(n)=\dsum_{i=1}^{n}(2 i - 1) = n^2$, which is defining $P(n)$ to be $\dsum_{i=1}^{n}(2 i - 1)$
and saying that it is also equal to $n^2$.  These are {\bf not} saying the same thing.
With the quotes, $P(n)$ is a propositional function.  Without them, it is a function from
$\BBZ$ to $\BBZ$.

In fact, to avoid this confusion, I recommend that you never use the equals sign with propositional
functions, especially when writing induction proofs.
\end{rem}

Now it's your turn to try to fill in the details of an induction proof.

\newpage % On own page!

\begin{fib}
Reprove Theorem~\ref{thm:sum-of-first-n-integ} using induction.  That is, prove that
for $n\geq 1$, $\dsum_{i=1}^{n} i = \dfrac{n(n+1)}{2}.$
\begin{pf}
Let $P(k)$ be the statement ``$\dsum_{i=1}^{k} i = \dfrac{k(k+1)}{2}$''.
We need to show that $P(n)$ is true for all $n\geq 1$.

\noindent
{\bf Base Case:} 
When $k=1$, we have
$\dsum_{i=1}^1 i = 1 = \blankline{1}$.  Therefore, \blankline{2}.

\noindent
{\bf Inductive Hypothesis:} 
Let $k\geq 1$, and assume that  \blankline{1.75}.
That is, assume that \blankline{3}.

\noindent
\begin{quote}
${[}$This is not part of the proof, but it will help us see what's next.
Our goal in the next step is to prove that 
\blankline{.75} is true.  That is, we need to show that
\blankline{2.5}.${]}$
\end{quote}
{\bf Inductive Step:}
Notice that
\begin{eqnarray*}
\dsum_{i=1}^{k+1}i&=&\blankline{1.5} + (k+1)\\
&=&\blankline{2}+(k+1) \mbox{(by the inductive hypothesis)}\\
&=&(k+1)\underline{\left(\hspace{1.5in}\rule{.000001in}{.35in}\right)}\\
&=&\blankline{2}
\end{eqnarray*}
Thus, \blankline{1.5}.

\noindent
{\bf Summary: }
We showed that \blankline{1.25} and that whenever \blankline{.75},
$P(k)\rightarrow P(k+1)$, 
therefore $P(n)$ is true for \blankline{1} by
\blankline{1.25}
\end{pf}
\begin{answer}
$\dfrac{1(1+1)}{2}$;
$P(1)$ is true;
$P(k)$ is true;
$\dsum_{i=1}^{k} i = \dfrac{k(k+1)}{2}$;
$P(k+1)$;
$\dsum_{i=1}^{k+1} i = \dfrac{(k+1)(k+2)}{2}$;
$\dsum_{i=1}^{k} i$;
$\dfrac{k(k+1)}{2}$;
$\dfrac{k}{2}+1$;
$\dfrac{(k+1)(k+2)}{2}$;
$P(k+1)$ is true;
$P(1)$ is true;
$k\geq 1$;
all $n\geq 1$;
induction {\em or} the principle of mathematical induction {\em or} PMI.
\end{answer}
\end{fib}

\subsection{Equalities/Inequalities}
The last few example induction proofs have dealt with statements of the form
$$LHS(k)=RHS(k),$$
where $LHS$ stands for {\em left hand side} 
and $RHS$ stands for {\em right hand side}.
For instance, in Example~\ref{exa:sumodd},
the statement was
$$\dsum_{i=1}^{n}(2 i - 1) = n^2,$$
so $LHS(k)=\dsum_{i=1}^{k}(2 i - 1)$ and $RHS(k) = k^2$.


\begin{ques}
Let $P(n)$ be the statement 
``$\dsum_{i=1}^{n} i\cdot i! = (n+1)!-1$.''
Determine each of the following:
\begin{lenum}
\item $P(k)$ is the statement \blankline{3}.
\item $P(k+1)$ is the statement \blankline{3}.
\item $LHS(k)=\blankline{3}$
\item $RHS(k)=\blankline{3}$
\item $LHS(k+1)=\blankline{3}$
\item $RHS(k+1)=\blankline{3}$
\end{lenum}
\begin{answer}
\begin{lenum}
\item $P(k)$ is the statement ``$\dsum_{i=1}^{k} i\cdot i! = (k+1)!-1$''
\item $P(k+1)$ is the statement ``$\dsum_{i=1}^{k+1} i\cdot i! = (k+2)!-1$''
\item $LHS(k)=\dsum_{i=1}^{k} i\cdot i!$
\item $RHS(k)=(k+1)!-1$
\item $LHS(k+1)=\dsum_{i=1}^{k+1} i\cdot i!$
\item $RHS(k+1)=(k+2)!-1$
\end{lenum}
\end{answer}
\end{ques}

For statements of this form, the goal of the inductive step is to show that 
$LHS(k+1)=RHS(k+1)$ given the fact that $LHS(k)$=$RHS(k)$ (the inductive hypothesis). 
The way this should generally be done is as follows:
\begin{proc}
Given a proposition of the form ``$LHS(n)=RHS(n)$,'' the algebra in the
inductive step of an induction proof should be done as follows:

\medskip
\noindent
\begin{tabular}{lllp{3.4in}}
$LHS(k+1)$&=&$LHS(k)+stuff$&(apply algebra to separate $LHS(k)$ from the rest)\\
        &=&$RHS(k)+stuff$&(use the inductive hypothesis to replace $LHS(k)$ with $RHS(k)$)\\
        &=&$\cdots$&(1 or more steps, usually involving algebra, that\\
        &=&$RHS(k+1)$&result in the goal of getting to $RHS(k+1)$)
\end{tabular}
\end{proc}
The last few examples followed this procedure, and your proofs should also follow it.
Notice that these examples {\em do not} begin the inductive step by writing out $LHS(k+1)=RHS(k+1)$.
One of them wrote it out, but it was {\em before} the inductive step for the purpose of
making the goal in the inductive step clear. 
The inductive step should always begin by writing just 
$LHS(k+1)$, and should then use algebra, the inductive hypothesis, etc., 
until $RHS(k+1)$ is obtained. 

This technique also works (with the appropriate slight modifications) 
with inequalities, e.g.
$$LHS(k)\leq RHS(k) \mbox{ and}$$
$$LHS(k)\geq RHS(k).$$ 
For instance, if $P(k)$ is the statement ``$k>2^k$'', $LHS(k)=k$, and $RHS(k)=2^k$.
In addition, the `$+stuff$' is not always literally addition.  For instance, it might be
$LHS(k)\times stuff$.

Here is another example of this type of induction proof--this time using an 
inequality. 

\begin{exa}
Prove that $n<2^n$ for all integers $n\geq 1$.
\begin{pf} 
Let $P(n)$ be the statement ``$n<2^n$''.
We want to prove that $P(n)$ is true for all $n\geq 1$.

\noindent
{\bf Base Case:} Since $1<2^1$, $P(1)$ is clearly true.

\noindent
{\bf Hypothesis:} We assume $P(k)$ is true if $k\geq 1$.  That is, $k<2^k$.
\begin{quote}
{\em Next we need to show that $P(k+1)$ is true.
That is, we need to show that $(k+1)<2^{k+1}$.  (Notice that I did not
state that this was true, and I do not start with this statement in the next step.
I am merely pointing out what I {\em need to prove}.)  This paragraph is not
really part of the proof--think of it as a side-comment or scratch work.}
\end{quote}
{\bf Inductive:}
Given that $k<2^k$, we can see that 
$$
\begin{array}{llll}
k+1&<&2^k+1&\mbox{(since $k<2^k$)}\\
   &<&2^k+2^k&\mbox{(since $1<2^k$ when $k\geq 1$)}\\
   &=&2(2^k)&\\
   &=&2^{k+1}&
\end{array}
$$
Since we have shown that $k+1<2^{k+1}$, $P(k+1)$ is true.

\noindent
{\bf Summary:} 
Since we proved that $P(1)$ is true, and that
$P(k)\rightarrow P(k+1)$,
by {\em PMI}, $P(n)$ is
true for all $n\geq 1$.
\end{pf}
\end{exa}

In the previous example, $LHS(k)=k$, so $LHS(k+1)$ is already in the form
$LHS(k)+stuff$ since $LHS(k+1)=k+1=LHS(k)+1$. So the first step of algebra is unnecessary
and we were able to apply the inductive hypothesis immediately.
Don't let this confuse you.  This is essentially the same as the other examples minus 
the need for algebra in the first step.

\begin{rem}
By the time you are done with this section, you will likely be tired of hearing this, 
but since it is the most common mistake made in induction proofs, it is worth repeating {\em ad nauseam}.
{\bf Never begin the inductive step of an induction proof by writing down $\mathbf{P(k+1)}$.}
You do not know it is true yet, so it is not valid to write it down as if it were true so that
you can use a technique such as working both sides to verify that it is true (which, as we have
also previously stated, is not a valid proof technique).

You {\bf can} (and sometimes {\bf should}) write down $P(k+1)$ on another piece of paper
or with a comment such as ``We need to prove that'' preceding it
so that you have a clear direction for the inductive step.
\end{rem}

If you can complete the next exercise without too much difficulty, 
you are well on your way to understanding how to write induction proofs.

%\newpage

\begin{exer}
Use induction to prove that for all $n\geq 1$, 
$\dsum _{i=1}^n i^2 = \dfrac{n(n+1)(2n+1)}{6}$.

\noindent
(Hint: Follow the techniques and format of the previous examples and 
be smart about your algebra and it will go a lot easier.
Also, you will need to factor a polynomial in the inductive step, 
but if you determine what the goal is ahead of time, it shouldn't
be too difficult.)

\vspace*{6.5in}
\begin{answer}
Define $P(n)$ to be the statement ``$\dsum _{i=1}^n i^2 = \dfrac{n(n+1)(2n+1)}{6}$''.
We need to show that $P(n)$ is true for all $n\geq 1$.

\noindent
{\bf Base Case:} 
Since $\dsum _{i=1}^1 i^2 = 1^2 = 1 =  \dfrac{1(2)(3)}{6}$,
$P(1)$ is true. (If your algebra is in a different order, like 
$\dsum _{i=1}^1 i^2 = \dfrac{1(2)(3)}{6}= 1$, it is incorrect.  We only know that 
$\dsum _{i=1}^1 i^2 = \dfrac{1(2)(3)}{6}$ because we first saw that $\dsum _{i=1}^1 i^2=1$,
and then were able to see that $1 =  \dfrac{1(2)(3)}{6}$.)

\noindent
{\bf Inductive Hypothesis:}
Let $k\geq 1$ and assume $P(k)$ is true.  That is, $\dsum _{i=1}^k i^2 = \dfrac{k(k+1)(2k+1)}{6}$.

\noindent
(As a side note, I know that what I need to prove next is

\noindent
$\displaystyle\dsum _{i=1}^{k+1} i^2 = \dfrac{(k+1)(k+2)(2(k+1)+1)}{6}=\dfrac{(k+1)(k+2)(2k+3)}{6}.$

\noindent
I am only writing this down now so that I know what my goal is.  I am not going to start working
both sides of this or otherwise manipulate it.  I can't because I don't know whether or not it is
true yet.)

\noindent
{\bf Inductive Step:}
Notice that 
\begin{eqnarray*}
\dsum _{i=1}^{k+1} i^2&=&\dsum _{i=1}^k i^2+(k+1)^2\\
&=&\dfrac{k(k+1)(2k+1)}{6}+(k+1)^2\\
&=&(k+1)\left(\dfrac{k(2k+1)}{6}+(k+1)\right)\\
&=&(k+1)\left(\dfrac{k(2k+1)+6(k+1)}{6}\right)\\
&=&(k+1)\left(\dfrac{2k^2 + k + 6k+ 6}{6}\right)\\
&=&(k+1)\left(\dfrac{2k^2 + 7k+ 6}{6}\right)\\
&=&(k+1)\left(\dfrac{(2k+3)(k+2)}{6}\right)\\
&=&\dfrac{(k+1)(k+2)(2k+3)}{6}.
\end{eqnarray*}
Therefore $P(k+1)$ is true.

\noindent
{\bf Summary:} Since $P(1)$ is true and $P(k)\rightarrow P(k+1)$ is true when $k\geq 1$,
$P(n)$ is true for all $n\geq 1$ by induction.
\end{answer}
\end{exer}

\subsection{Variations}

In this section we will discuss a few slight variations of the details we have presented so far.
First we discuss the fact that we do not need to use a propositional function.
Then we will discuss a variation regarding the inductive hypothesis.

It is not always necessary to explicitly define $P(k)$ for use in an induction proof.
$P(k)$ is used mostly for convenience and clarity. For instance, in the solution to the previous
exercise, it allowed us to just say 
\begin{center}
``$P(k)$ \mbox{is true}''
\end{center}
instead of saying 
\begin{center}
``$\dsum _{i=1}^n i^2 = \dfrac{n(n+1)(2n+1)}{6}$''\, \,  (which is long)
\end{center}
or
\begin{center}
``the statement is true for $k$''\, \,  (which is a little vague/awkward).
\end{center}

Here is an example that does not use $P(k)$.  It also does not label the four parts of the proof.  
That is perfectly fine.  The main reason we have done so in previous examples 
is to help you identify them more clearly.

\begin{exa}\index{Fibonacci numbers}
Let $f_n$ be the $n$-th Fibonacci number.
Prove that for all integers $n \geq 1,$
$$f_{n - 1}f_{n + 1} = f_n ^2 + (-1)^{n}.$$ 
\begin{pf}
For $k = 1,$ we have
$$f_0f_2= 0\cdot 1 = 0 = 1 - 1 = 1^2 + (-1)^1 = f_1^2 +  (-1)^1,$$
and so the assertion is true for $k = 1.$ 
Suppose $k \geq 1,$ and  that the
assertion is true for $k$.  That is,
$$f_{k - 1}f_{k + 1} = f_k ^2 + (-1)^{k}.$$
This can be rewritten as 
$$f_k ^2  = f_{k - 1}f_{k + 1} -(-1)^{k}$$
(a fact that we will find useful below).
Then 
$$\begin{array}{llll}
f_{k}f_{k + 2} & =  & f_k(f_{k + 1} + f_k) &\mbox{(by definition of $f_n$ applied to $f_{k+2}$)}\\
& = & f_kf_{k + 1} + f_k ^2\\
& = & f_kf_{k + 1} + f_{k - 1}f_{k + 1} - (-1)^{k} &\mbox{(by rewritten inductive hypothesis)}\\
& = & f_{k + 1}(f_k + f_{k - 1}) + (-1)^{k + 1} \\
& = & f_{k + 1}f_{k + 1} + (-1)^{k + 1}&\mbox{(by the definition of $f_k$)}\\
& = & f_{k + 1}^2 + (-1)^{k + 1},
\end{array}$$
and so the assertion is true for $k+1$.
The result follows by induction.
\end{pf}
\end{exa}

\begin{exer}\label{exer:sumi2i}
Use induction to prove  that for all $n\geq 1$, 
$$1\cdot 2 + 2\cdot 2^2 + 3\cdot 2^3 + \cdots + n\cdot 2^n = 2 + (n - 1)2^{n + 1}$$
or if you prefer, 
$$\sum_{i=1}^n i\cdot2^i =  2 + (n - 1)2^{n + 1}.$$
Do so without using a propositional function.
You may label the four parts of your proof, but it is not required.

\vspace*{6.5in}
\begin{answer}
For $k = 1$ we have $1\cdot 2 = 2 + (1-1)2^2$, and so the
statement is true for $n = 1$. 
Let $k\geq 1$ and assume the statement is true for
$k$. That is, assume
$$1\cdot 2 + 2\cdot 2^2 + 3\cdot 2^3 + \cdots + k\cdot 2^k = 2 + (k - 1)2^{k + 1}.$$
We need to show that
$$1\cdot 2 + 2\cdot 2^2 + 3\cdot 2^3 + \cdots + (k+1)\cdot 2^{k+1} = 2 + k2^{k + 2}.$$
Using some algebra and the inductive hypothesis, we can see that 
\begin{eqnarray*}
1\cdot 2 + 2\cdot 2^2 + 3\cdot 2^3 + \cdots + k\cdot 2^k + (k+1)2^{k+1} 
&=& 2 + (k - 1)2^{k + 1} + (k+1)2^{k+1}\\
&=& 2 + (k - 1+k+1)2^{k + 1}\\
& =& 2 + 2k2^{k+1}\\ 
&=& 2 + k2^{k+2}.
\end{eqnarray*}
Thus, the result is true for $k+1$.  The result follows by induction.
\end{answer}
\end{exer}


\begin{exa}
Prove the generalized form of DeMorgan's law.  That is, show that
for any $n\geq 2$, if $p_1$, $p_2$, $\ldots$, $p_n$ are propositions, then 
$$\neg(p_1\vee p_2 \vee\cdots\vee p_n)=(\neg p_1\wedge \neg p_2 \wedge \cdots\wedge\neg p_n).$$

We provide several appropriate proofs of this one (and one inappropriate one).

\medskip
\noindent
{\bf Proof 1:} (A typical proof)
\begin{quote}
Let $P(n)$ be the statement
``$\neg(p_1\vee p_2 \vee\cdots\vee p_n)=(\neg p_1\wedge \neg p_2 \wedge \cdots\wedge\neg p_n)$.''
We want to show that for all $n\geq 2$, $P(n)$ is true.
$P(2)$ is DeMorgan's law, so the base case is true.
Assume $P(k)$ is true.  Then 
$$
\begin{array}{llll}
\neg(p_1\vee p_2 \vee\cdots\vee p_{k+1})
&=&
\neg((p_1\vee p_2 \vee\cdots\vee p_{k})\vee p_{k+1})
&\mbox{associative law}\\
&=&
\neg(p_1\vee p_2 \vee\cdots\vee p_{k})\wedge \neg p_{k+1}
&\mbox{DeMorgan's law}\\
&=&
(\neg p_1\wedge \neg p_2 \wedge \cdots\wedge\neg p_{k})\wedge \neg p_{k+1}
&\mbox{hypothesis}\\
&=&
(\neg p_1\wedge \neg p_2 \wedge \cdots\wedge\neg p_{k}\wedge \neg p_{k+1})
&\mbox{associative law}
\end{array}
$$
Thus $P(k+1)$ is true.
Since we proved that $P(2)$ is true, and that
$P(k)\rightarrow P(k+1)$ if $k\geq 2$,
by {\em PMI}, $P(n)$ is true for all $n\geq 2$.\hfill$\square$
\end{quote}

\noindent
{\bf Proof 2:} (Not explicitly defining/using $P(n)$)
\begin{quote}
We know that 
$\neg(p_1\vee p_2)=(\neg p_1\wedge \neg p_2)$ since
this is simply DeMorgan's law.
Assume the statement is true for $k$.  That is, 
$\neg(p_1\vee p_2 \vee\cdots\vee p_{k})=
(\neg p_1\wedge \neg p_2 \wedge \cdots\wedge\neg p_{k})$.
Then we can see that 
$$
\begin{array}{llll}
\neg(p_1\vee p_2 \vee\cdots\vee p_{k+1})
&=&
\neg((p_1\vee p_2 \vee\cdots\vee p_{k})\vee p_{k+1})
&\mbox{associative law}\\
&=&
\neg(p_1\vee p_2 \vee\cdots\vee p_{k})\wedge \neg p_{k+1}
&\mbox{DeMorgan's law}\\
&=&
(\neg p_1\wedge \neg p_2 \wedge \cdots\wedge\neg p_{k})\wedge \neg p_{k+1}
&\mbox{hypothesis}\\
&=&
(\neg p_1\wedge \neg p_2 \wedge \cdots\wedge\neg p_{k}\wedge \neg p_{k+1})
&\mbox{associative law}\\
\end{array}
$$
Thus the statement is true for $k+1$. Since we have shown that the
statement is true for $n=2$, and that whenever it is true for $k$ it
is true for $k+1$, 
by {\em PMI}, the statement is true for all $n\geq 2$.\hfill$\square$

{\em Sometimes it is acceptable to omit the justification in the summary.
That is, there isn't necessarily a need to restate what you have proven and
you can just jump to the conclusion.
So the previous proof could end as follows:}

\begin{quote}
Thus the statement is true for $k+1$. 
By {\em PMI}, the statement is true for all $n\geq 2$.
\end{quote}
\end{quote}

%\noindent
%{\bf Proof 3:} (Algebra steps not justified)
%\begin{quote}
%The case $k=2$ is DeMorgan's law, and if the statement is true for $k$, then
%{\small $$
%\neg(p_1\vee\cdots\vee p_{k+1}) =
%\neg(p_1\vee \cdots\vee p_{k})\wedge \neg p_{k+1} =
%(\neg p_1\wedge \cdots\wedge\neg p_{k})\wedge \neg p_{k+1} =
%(\neg p_1\wedge \neg p_2 \wedge \cdots\wedge\neg p_{k}\wedge \neg p_{k+1})
%$$
%}
%Thus the statement is true for $k+1$. 
%By {\em PMI}, the statement is true for all $n\geq 2$.\hfill$\square$
%\end{quote}

\noindent 
{\bf Proof 3:} (common in journal articles, unacceptable for this class)
\begin{quote}
The result follows easily by induction.\hfill$\square$
\end{quote}
\end{exa}

\newpage % Force it because it looks crappy otherwise.

\begin{eval}
Prove that for all positive integers $n$,
${\displaystyle \sum_{i=1}^{n} i\cdot i! = (n+1)!-1}$.
\begin{ssol}
Base: $n=1$
\begin{eqnarray*}
1\cdot 1!&=&(1+1)!-1\\
1&=&2!-1\\
1&=&1
\end{eqnarray*}
Assume $\dsum_{i=1}^n i\cdot i! = (n+1)!-1$ for $n\geq 1$.\\
Induction:
\begin{eqnarray*}
\sum_{i=1}^{n+1}i\cdot i! &=& \sum_{i=1}^{n}i\cdot i! + (n+1)(n+1)!\\
&=& (n+1)! - 1 + (n+1)(n+1)!\\
&=&  (n+1+1)(n+1)! -1\\
&=&  (n+2)(n+1)! -1\\
&=&  (n+2)! -1
\end{eqnarray*}
Therefore it is true for $n$.  Thus by PMI it is true for $n\geq 1$.
\end{ssol}
\evallinethree
\begin{answer}
This proof is very close to being correct, but is suffers from a few small but important errors:
\begin{itemize}
\item
For the sake of clarity, it might have been better to use $k$ throughout most of the proof
instead of $n$.  The exception is in the final sentence where $n$ is correct.
\item
The base case is just some algebra without context.  A few words are needed.  For instance, 
`notice that when $n=1$,'.
\item
The base case is presented incorrectly.  Notice that the writer starts by writing down what
she wants to be true and then deduces that it is indeed correct by doing algebra on both sides
of the equation.  As we have already mentioned, 
{\em you should {\bf never} start with what you want to prove and work both sides!}
It is not only sloppy, but it can lead to incorrect proofs.
Whenever I see students do this, I always tell them to use what I call the $U$ method.  What I mean
is rewrite your work by starting at the upper left, going down the left side, then doing up the
right side.  So the above should be rewritten as:
$$1\cdot 1! = 1 = 2!-1 = (1+1)!-1.$$
Notice that if the $U$ method does not work (because one
or more steps isn't correct), it is probably an indication
of an incorrect proof.  Consider what happens if you try it
on the proof in Exercise~\ref{minusOneIsNotOne}.  You would
write $-1=(-1)^2=1=1^2=1$.  Notice that the first equality
is incorrect.

The $U$ method can sometimes apply to inequalities as well.
\item
When the writer makes her assumption, she says `for $n\geq 1$'.  This is O.K., but there is some ambiguity
here.  Does she mean for {\em all} $n$, or for a particular value of $n$?  She must mean the latter since
the former is what she is trying to prove.  It would have been better for her to say `for some $n\geq 1$.'
\item
The algebra in the inductive step is perfect.  However, what does it mean?  She should include something like
`Notice that' before her algebra just to give it a little context.  It often doesn't take a lot of words,
but adding a few phrases here and there goes a long way to help a proof flow more clearly.
\item
She says `Therefore it is true for $n$'.  She must have meant $n+1$ since that is what she just proved.
\item
As with her assumption, her final statement could be clarified by saying `for all $n\geq 1$.'
\end{itemize}
Overall, the proof has almost all of the correct content.  
Most of the problems have to do with presentation.
But as we have seen with
other types of proofs, 
the details are really important to get right!
\end{answer}
\end{eval}

\vspace*{.125in}
The second variation we wish to discuss has to do with the inductive hypothesis/step.
In the inductive step, we can replace 
$P(k)\rightarrow P(k+1)$ with 
$P(k-1)\rightarrow P(k)$ as long as we prove the statement
for all $k$ {\em larger than any of the base cases}.
In general, we can use whatever index we want for the inductive hypothesis as long as we use it to prove that the statement
is true for the next index, and as long as we are sure to
cover all of the indices down to the base case.
For instance, if we prove $P(k+3)\rightarrow P(k+4)$, then
we need to show it for all $k+3\geq a$ (that is, all $k\geq a-3$), assuming $a$ is the base case.
Put simply, the assumption we make about the value of $k$
must guarantee that the inductive hypothesis includes the base case(s).



%Make sure the value of $k$ makes the hypothesis thingy include the base case.  haha.  lol.  jk.

\begin{ques}
Consider a `proof' of $\forall n P(n)$
that shows that $P(1)$ is true and 
that $P(k)\rightarrow P(k+1)$ for $k>1$.
What is wrong with such a proof?

\noindent
\answerlinetwo
\begin{answer}
Given this proof, we know that $P(1)$ is true.
We also know that 
$P(2)\rightarrow P(3)$, $P(3)\rightarrow P(4)$, etc, 
are all true.
Unfortunately, the proof omits showing that $P(2)$ is true,
so modus ponens never applies.   In other words, knowing that
$P(2)\rightarrow P(3)$ is true does us no good unless we know
$P(2)$ is true, which we don't.  Because of this, we don't
know anything about the truth values of $P(3)$, $P(4)$, etc.
The proof either needs to show that $P(2)$ is true as part of
the base case, or the inductive step needs to start at $1$
instead of $2$. 
\end{answer}
\end{ques}

\begin{rem}
Whether you assume $P(k)$ or $P(k-1)$ is true, 
you must specify the values of $k$ precisely based on your choice.  
For instance, if you assume $P(k)$ is true for all $k>a$, you have a problem.
Although you known $P(a)$ is true (because it is a base case), when you assume 
$P(k)$ is true for $k>a$, the smallest $k$ can be is $a+1$.  In other words, 
when you prove $P(k)\rightarrow P(k+1)$, you leave out $P(a)\rightarrow P(a+1)$.
But that means you can't get anywhere from the base case, so the whole proof is invalid.
\end{rem}

If you are wondering why we would use $P(k-1)$ as the inductive
hypothesis instead of $P(k)$, it is because sometimes it makes
the proof easier--for instance, the algebra steps involved might 
be simpler.  


\begin{exa} 
Prove that the expression $$ 3^{3n + 3} - 26n -
27$$is a multiple of $169$ for all natural numbers $n$.
\begin{pf}
Let $P(k)$ be the statement ``$3^{3k + 3} - 26k - 27 = 169 N$ for some $ N\in\BBN$.'' 
We will prove that $P(0)$ is true and that 
$P(k - 1) \rightarrow P(k)$.

\noindent
%When $k = 1$ notice that $3^{3\cdot 1 + 3} - 26\cdot 1 - 27= 676 = 169\cdot 4$, so $P(1)$ is true.

\noindent
When $k = 0$, $3^{3\cdot 0 + 3} - 26\cdot 0 - 27= 27-27=0 = 169\cdot 0$, so $P(0)$ is true.

Let $k> 0$ and assume $P(k - 1)$ is true.
That is, there is some $N\in\BBN$ such that 
$3^{3(k - 1) + 3} - 26(k - 1) - 27 = 169N$.  After a little algebra, this is
the same as $3^{3k} - 26k - 1 = 169N$. Then
\begin{eqnarray*}
3^{3k + 3} - 26k - 27 
&=& 27\cdot 3^{3k} - 26k - 27\\
&=& 27\cdot 3^{3k} +(26-27)26k - 27\\
&=& 27\cdot 3^{3k} -27\cdot26 k - 27 + 26\cdot 26k\\
& =& 27(3^{3k} - 26k -1) + 676k\\
&=& 27\cdot 169N + 169\cdot 4k \, \, \, \mbox{(By the inductive hypothesis)}\\
&=& 169(27\cdot N +\cdot 4k)
\end{eqnarray*}
which is divisible by 169. The assertion is thus established by
induction.
\end{pf}
\end{exa}

\begin{ques}
Did you notice that in the previous example we assumed $k>0$ instead of $k\geq 0$?
Why did we do that?

\noindent
\answerline
\begin{answer}
Because our inductive hypothesis was that $P(k-1)$ is true instead of $P(k)$.  
If we assumed that $k\geq 0$, then when $k=0$ it would mean we are assuming $P(-1)$ is true, and
we don't know whether or not it is since we never discussed $P(-1)$.
\end{answer}
\end{ques}


\subsection{Strong Induction}
The form of induction we have discussed up to this point only assumes the statement is true for 
one value of $k$.  This is sometimes called {\em weak induction}.\index{weak induction}
In {\em strong induction}, we assume that the statement is true for all values up to and including $k$.
\index{strong induction}
In other words, with strong induction, the inductive hypothesis involves proving that 
$$[P(a)\wedge P(a+1)\wedge \cdots \wedge P(k)]\rightarrow P(k+1) 
\mbox{ if } k\geq a.$$
This may look more complicated, but practically speaking, there is really very little difference. 
Essentially, strong induction just allows us to assume {\em more} than weak induction.  
Let's see an example of why we might need strong induction.

\begin{exa}
Show that every integer $n\geq 2$ can be written as the product of primes.
\begin{pf}
Let $P(n)$ be the statement ``$n$ can be written as the product of primes.''
We need to show that for all $n\geq 2$, $P(n)$ is true.

Since 2 is clearly prime, it can be written as the product of one
prime.  Thus $P(2)$ is true.

Assume $[P(2)\wedge P(3)\wedge\cdots \wedge P(k-1)]$ is true for $k>2$.
In other words, assume all of the numbers from $2$ to $k-1$ can be written as the product of primes.

We need to show that $P(k)$ is true. 
If $k$ is prime, clearly $P(k)$ is true.  If $k$ is not prime,
then we can write $k=a\cdot b$, where $2\leq a\leq b< k.$
By hypothesis, $P(a)$ and $P(b)$ are true, so $a$ and $b$ can be written as the product of primes.
Therefore, $k$ can be written as the product of primes, namely the primes from the
factorizations of $a$ and $b$. Thus $P(k)$ is true.

Since we proved that $P(2)$ is true, and that
$[P(2)\wedge P(3)\wedge\cdots\wedge P(k-1)]\rightarrow P(k)$ if $k> 2$,
by the principle of mathematical induction, $P(n)$ is true for all $n\geq 2$.
That is, every integers $n\geq 2$ can be written as the product of primes.
\end{pf}
\end{exa}


\begin{exa}
In the country of SmallPesia coins only come in values of $3$ and
$5$ pesos. Show that any quantity of pesos greater than or equal
to $8$ can be paid using the available coins.
\begin{pf}
{\bf Base Case:} Observe that $8 = 3 + 5$, $9 = 3 + 3 + 3$, and $10 = 5 + 5$, 
so we can pay $8$, $9$, or $10$ pesos with the available coinage.

{\bf Inductive Hypothesis:} Assume we can pay any value from $8$ to $k-1$ pesos, where
$k\geq 11$.  

{\bf Inductive step:}
The inductive hypothesis implies that we can pay with $k-3$ pesos.
We can add to the coins used for $k-3$ pesos a single coin of value
$3$ in order to pay for $k$ pesos.

{\bf Summary:}
Since we can pay for $8$, $9$, and $10$ pesos, and whenever we can pay
for anything between $8$ and $k-1$ pesos we can pay for $k$ pesos, the
strong form of induction implies that we can pay for any quantity of
pesos $n\geq 8$.

{\em 
Notice that the reason we needed three base cases for this proof was
the fact that we looked back at $k-3$, three value previous to the value of
interest.  If we had only proven it for $8$, we would have needed to
prove $9$ and (more importantly) $10$ in the inductive step.  
But the inductive step doesn't work for $10$ since there is no solution for $10-3=7$ pesos.
}
\end{pf}
\end{exa}

Notice that there is no way we could have used weak induction in either of the previous examples.

\newpage
\subsection{Induction Errors}

The following examples should help you appreciate why we need 
to be very precise/careful when writing induction proofs.



\begin{exa}
What is wrong with the following (supposed) proof that $a^n=1$ for $n\geq 0$:
\begin{pf}
{\em Base case:} Since $a^0=1$, the statement is true for $n=0$.

\noindent
{\em Inductive step:} Let $k>0$ and assume $a^j=1$ for $0\leq j\leq k$.
Then
$$a^{k+1} = \frac{a^k\cdot a^k}{a^{k-1}}=\frac{1\cdot 1}{1} = 1.$$
{\em Summary:} Therefore by PMI, $a^n=1$ for all $n\geq 0$.
\end{pf}
\begin{sol}
The base case is correct, and there is nothing wrong with the summary, assuming the
inductive step is correct.
$a^k=1$ and $a^{k-1}=1$ are correct by the inductive hypothesis
since we are assuming $k>0$.  The algebra is also correct.
So what is wrong?
The problem is that when  $k=0$, $a^{-1}$ would be in the denominator.
But we don't know whether or not $a^{-1}=1$.  Thus we needed to assume $k>0$.
As it turns out, that is precisely where the problem lies.
We proved that $P(0)$ is true and that $P(k)\rightarrow P(k+1)$ is true when $k>0$.
Thus, we know that $P(1)\rightarrow P(2)$, and $P(2)\rightarrow P(3)$, etc., but 
we never showed that $P(0)\rightarrow P(1)$ because, of course, it isn't true.
The induction doesn't work without $P(0)\rightarrow P(1)$.
\end{sol}
\end{exa}


\begin{eval}
Prove or disprove that all goats are the same color.
\begin{ssol}
If there is one goat, it is obviously the same color as itself.
Let $n\geq 1$ and assume that any collection of $n$ goats are
all the same color.  Consider a collection of $n+1$ goats.
Number the goats $1$ through $n+1$.
Then goats $1$ through $n$ are the same color (since there
are $n$ of them) and goats $2$ through $n+1$ are the same
color (again, since there are $n$ of them).
Since goat 2 is in both collections, the goats in both collections
are the same color.  Thus, all $n+1$ goats are the same color.
\end{ssol}
\evallinethree
\begin{answer}
This contains a very subtle error.  Did you find it?
If not, go back and carefully re-read the proof and think
carefully--at least one thing said in the proof {\em must}
be incorrect.  What is it?

O.K., here it is:  The statement `goat 2 is in both collections'
is not always true.  If $n=1$, then the first collection
contains goats $1$ through $1$, and the second collection
contains goats $2$ through $2$.  In this case, there is no overlap
of goat $2$, so the proof falls apart.
\end{answer}
\end{eval}



The next example deals with {\em binary palindromes}.
Binary palindromes can be defined recursively by $\lambda, 0, 1\in P$, and whenever $p\in P$, then
$1p1\in P$ and $0p0\in P$.  
(Note: $\lambda$ is the notation sometimes used to denote the {\em empty string}---that is, the string of length $0$.  
Also, $1p1$ means the binary string obtained by appending $1$ to the begin and end of string $p$.  Similarly for $0p0$.)
Notice that there is 1 palindrome of length 0 ($\lambda$), 2 of length 1 (0, 1),
2 of length 2 (00, 11), 4 of length 3 (000, 010, 101, 111), etc.

\begin{eval}
Use induction to prove that the number of binary palindromes of length $2n$ (even length) is 
$2^n$ for all $n\geq 0$.
\begin{spfenum}
\item
Base case: $k=1$.  The total number of palindromes of length $2=2$ is $2^1=2$.  It is true.

\noindent
Assume the total number of binary palindromes with length $2k$ is $2^k$.  To form a binary palindrome
with length $2(k+1)=2k+2$, with every element in the set of binary palindromes with length $2k$ we either
put $(00)$ or $(11)$ to the end or beginning of it.  
Therefore, the number of binary palindromes with length $2(k+1)$ is twice as many as the number of
binary palindromes with length $2k$, which is $2\times2^k=2^{k+1}$.  Thus it is true for $k+1$.
By the principle of mathematical induction, the total number of binary palindromes of length
$2n$ for $n\geq 1$ is $2^n$.

\noindent
\evallinethree
\item
For the base case, notice that there is $1=2^0$ palindromes of length $0$ (the empty string).
Now assume it is true for all $n$.
For each consecutive binary number with $n$ bits, you are adding a bit to either end, which multiplies
the total number by $2^2$ permutations, but for it to be a palindrome, they both have to be either
$0$ or $1$, so it would just be $2$ instead, so for binary numbers of length $2k$, there are $2^k$ 
palindromes.

\noindent
\evallinethree
\item
The empty string is the only string of length $0$, and it is a palindrome.  
Thus there is $1=2^0$ palindromes of length $0$.
Let $2n$ be the length, assume $2n\rightarrow 2^n$ palindromes.
Now we look at $n+1$ so we know the length is $2n+2$ and it starts and ends with
either $0$ or $1$ and has $2n$ values in between.
Both possibilities imply $2^n$ palindromes, so $2^n+2^n=2^{n+1}$.

\noindent
\evallinethree
\end{spfenum}

\begin{answer}
\begin{pfevalenum}
\item
This solution is on the right track,  but it has several technical problems.  
	\begin{itemize}
	\item
	The base case should
	be $k=0$, not $k=1$.  
	\item
	The way the base case is worded could be improved.  For instance, 
	what purpose does saying `$2=2$' serve?  Also, the separate sentence that just says `it is true' is
	a little vague and awkward.  I would reword this as:
		\begin{quote}
		The total number of palindromes of length $2\cdot 1$ is $2=2^1$, so the statement is true for $k=1$.
		\end{quote}
	Of course, the base case should {\em really} be $k=0$, but if it were $k=1$, that is how I would word it.
	\item 
	The connection between palindromes of length $2k$ and $2(k+1)$ is not entirely clear and is incorrect
	as stated.  A palindrome of length $2(k+1)$ can be formed from a palindrome of length $2k$ by adding
	a $0$ to both the beginning and end or adding a $1$ to both the beginning and the end.  
	This what was probably meant, but it is not what the proof actually says. 
	 
	But we need to say a little more about this.  Every palindrome of length
	$2(k+1)$ can be formed from exactly one palindrome of length $2k$ with this method.  But is this enough?  
	Not quite.  We also need to know that every palindrome of length $2k$ can be extended to a palindrome of
	length $2(k+1)$, and it should be clear that this is the case.  
	In summary, the inductive step needs to establish that there are twice as many binary palindromes of length $2(k+1)$ as there are of length $2k$.  
	The argument has to convince the reader that there is a 2-to-1 correspondence between these sets of palindromes.  In other words, we did not omit or double-count any.
	\end{itemize}
\item
	The base case correct.  Unfortunately, that is about the only thing that is correct.
	\begin{itemize}
	\item 
	The second sentence is wrong. 
	We cannot say that `it is true for all $n$'--that is precisely what we are trying to prove.  We need
	to assume it is true for a {\em particular} $n$ and then prove it is true for $n+1$.
	\item
	The rest of the proof is one really long sentence that is difficult to follow.  It should be split
	into much shorter sentences, each of which provides one step of the proof.	
	\item
	The term `binary number' should be replaced with `binary palindrome' throughout.  It causes
	confusion, especially when the words `add' and `consecutive' are used.  These mean something
	very different if we have numbers in mind instead of strings.
	\item  
	I don't think the phrase `each consecutive binary number' means what the writer thinks it means. 
	The binary numbers $1001$ and $1010$ are consecutive (representing $9$ and $10$), but that is 
	probably {\em not}	what the writer has in mind.
	\item
	The term `permutations' shows up for some reason.  
	I think they might have mean `strings' or something else.
	\item
	Why bring up the $4$ possible ways to extend a binary
	string by adding to the beginning and end if only two of them are relevant? 
	Why not just consider the ones of interest in the first place?
	\item
	In the context of a proof, the phrase 'you are adding' doesn't make sense.  Why am I adding something
	and what am I adding it to?  And do they mean addition (of the binary numbers) or appending (of strings)?
	\item 
	They switch from $n$ to $k$ in the middle of the proof to provide further confusion.
	\end{itemize}
\item
This proof has most of the right ideas, but it does not put them together well.  
The base case is correct.
It sounds like the writer understands what is going on with the inductive step, 
but needs to communicate it more clearly.
More specifically, what does `assume $2k\rightarrow 2^k$ palindromes' mean?
I think I am supposed to read this as `assume that there are $2^k$ palindromes of length $2k$.'
\footnote{In general, avoid the use of mathematical symbols in constructing the grammar of an English sentence.
One of the most common abuses I see is the use of $\rightarrow$ in the middle of a sentence.}

The final sentence is also problematic.  The first phrase tries to connect to the previous sentence,
but the connection needs to be a little more clear.  The final phrase is not a complete thought.
In the first place, I know that  $2^k+2^k=2^{k+1}$ and this has nothing to do with the previous
phrases.  In other words, the `so' connecting the phrases doesn't make sense.  But more seriously,
why do I care that $2^k+2^k=2^{k+1}$?  What he meant was something like 
`so there are $2^k+2^k=2^{k+1}$ palindromes of length $2k+2$'.
\end{pfevalenum}
\end{answer}
\end{eval}

\newpage

\begin{exer}
Based on the feedback from the previous Evaluate exercise, construct a proper proof that 
the number of binary palindromes of length $2n$ is $2^n$ for all $n\geq 0$.

\vspace*{4.75in}
\begin{answer}
The empty string is the only string of length $0$, and it is a palindrome.  
Thus there is $1=2^0$ palindromes of length $0$.

Now assume there are $2^n$ binary palindromes of length $2n$.
For every palindrome of length $2n$, exactly two palindromes of length $2(n+1)$ can be constructed
by appending either a $0$ or a $1$ to both the beginning and the end.  Further, every palindrome
of length $2(n+1)$ can be constructed this way.  Thus, there are twice as many palindromes of
length $2(n+1)$ as there are of length $2n$.  By the inductive hypothesis, there are $2\cdot 2^n=2^{n+1}$
binary palindromes of length $2(n+1)$.

The result follows by PMI.
\end{answer}
\end{exer}

\subsection{Summary/Tips}

Induction proofs are both intuitive and non-intuitive.  On the one hand, when you talk through
the idea, it seems to make sense.  On the other hand, it almost seems like you are using
{\em circular reasoning}.  It is important to understand that induction proofs do {\em not}
rely on circular reasoning.  Circular reasoning is when you assume $p$ in order to prove $p$.
But here we are not doing that.  We are assuming $P(k)$ and using that fact to prove $P(k+1)$, 
a different statement.  However, we are {\em not} assuming that $P(k)$ is true for all $k\geq a$.  
We are proving that {\bf\em if we assume that $P(k)$ is true}, then $P(k+1)$ is true.  
The difference between these statements may seem subtle, but it is important.

Let's summarize our approach to writing an induction proof.  
This is similar to Procedure~\ref{proc:induct1} except we include several of the unofficial steps
we have been using that often come in handy.
You are not required to use this procedure, but if you are having a difficult time with
induction proofs, try this out.  Here is the brief version.  After this we provide some
further comments about each step.

\newpage % This should be together on one page.

\begin{proc}\label{proc:induct2}
A slightly longer approach to writing an induction proof is as follows:
\begin{enumerate}
\item{\bf Define:} (optional) 
Define $P(n)$ based on the statement you need to prove.
\item{\bf Rephrase:}  (optional)
Rephrase the statement you are trying to prove using $P(n)$. 
This step is mostly to help you be clear on what you need to prove.
\item{\bf Base Case:} 
Prove the base case or cases.
\item{\bf Inductive Hypothesis:} 
Write down the inductive hypothesis. Usually it is as simple as ``Assume that $P(k)$ is true''.
\item{\bf Goal:} (optional)
Write out the goal of the {\em inductive step} (coming next).
It is usually ``I need to show that $P(k+1)$ is true''
It can be helpful to explicitly write out $P(k+1)$, although see important comments
about this step below.
This is another step that is mostly for your own clarity.
\item{\bf Inductive:} Prove the {\em goal} statement, usually using the 
{\em inductive hypothesis}.
\item{\bf Summary:} The typical induction summary.
\end{enumerate}
\end{proc}


Here are some comments about the steps in Procedure~\ref{proc:induct2}.

\begin{enumerate}
\item{\bf Define:} 
$P(n)$ should be a statement
about a single instance, not about a series of instances.  For example,
it should be statements like 
``$2 n$ is even''  or 
``A set with $n$ elements has $2^n$ subsets.''
It should {\em NOT} be of the form 
``$2 n$ is even if $n>1$,'' ``$n^2>0$ if $n\neq 0$,''
or ``For all $n>1$, a set with $n$ elements has $2^n$ subsets.''
\item{\bf Rephrase:} 
In almost all cases, the rephrased statement should be 
``For all $n\geq a$, $P(n)$ is true,'' where $a$ is some constant,
often 0 or 1. If the statement cannot be phrased in this way, induction may
not be appropriate.
\item{\bf Base Case:} 
For most statements, this means showing that $P(a)$ is true, where $a$ is
the value from the rephrased statement.  
Although usually one base case suffices, sometimes one must prove multiple
base cases, usually $P(a)$, $P(a+1)$, \ldots, $P(a+i)$ for some $i>0$.
This depends on the details of the inductive step.
\item{\bf Inductive Hypothesis:}
This is almost always one of the following: 
\begin{itemize}
\item
Assume that $P(k) \mbox{ is true}$.
\item
Assume that $P(k-1) \mbox{ is true}$.
\item
Assume that $[P(a)\wedge P(a+1)\wedge \cdots \wedge P(k)] \mbox{ is true}$ (strong induction)
\end{itemize}
Sometimes it is helpful to write out the hypothesis explicitly (that is,
write down the whole statement with $k$ or $k-1$ plugged in).

\item{\bf Goal:}
As previously stated, this is almost always ``I need to show that $P(k+1)$ is true'' 
(or ``I need to show that $P(k)$ is true'').
But it can be very helpful to explicitly write out what $P(k+1)$ is so you have a clear
direction for the next step.  However, {\em it is very important that you do not just
write out $P(k+1)$ without prefacing it with a statement like ``I need to show that...''.}
Since you are about to prove that $P(k+1)$ is true, you don't know that it is true
yet, so writing it down as if it is a fact is incorrect and confusing.
In fact, it is probably better write the goal separate from the rest if the proof
(e.g. on another piece of paper).  

The goal does not need to be written down and is not really part of the proof.  The 
only purpose of doing so it to help you see what you need to do in the next step. 
For instance, knowing the goal often helps you to figure out the required algebra steps to get there.

%$$\mbox{I need to show that if } k\geq a, P(k)\rightarrow P(k+1)$$
%or
%$$\mbox{I need to show that if } k> a,[P(a)\wedge P(a+1)\wedge \cdots \wedge P(k)]\rightarrow P(k+1).$$
%As mentioned previously, this step can also be 
%$$\mbox{I need to show that if } k> a, P(k-1)\rightarrow P(k)$$
%if it makes the induction step easier (e.g. it simplifies the necessary algebra). 

\item{\bf Inductive:} 
This is the longest, and most varied, part of the proof.
Once you get the hang of induction, you will typically only think about two parts of the 
proof---the base case and this step.  The rest will become second nature.

The inductive step should {\em not} start with writing down $P(k+1)$.
Some students want to write out $P(k+1)$ and
work both sides until they get them to be the same.  
As we have emphasized on several occasions, this is {\em not} a proper
proof technique.  You cannot start with something you do not know and then work it until
you get to something you do know and then declare it is true.  
\item{\bf Summary:} 
This is easy.  It is almost always either:
\begin{quote}
``Since we proved that $P(a)$ is true, and that
$P(k)\rightarrow P(k+1),$ for $k\geq a$, 
then we know that $P(n)$ is true for all $n\geq a$ by {\em PMI}, '' or
\end{quote}
\begin{quote}
``Since we proved that $P(a)$ is true, and that
$[P(a)\wedge P(a+1)\wedge\cdots\wedge P(k)]\rightarrow P(k+1),$ for $k\geq a$, 
$P(n)$ is true for all $n\geq a$ by {\em PMI}.''
\end{quote}
The details change a bit depending on what your inductive hypothesis was
(e.g. if it was $P(k-1)$ instead of $P(k)$).
Technically speaking, you can just summarize your proof by saying
\begin{quote}
``Thus, $P(n)$ is true for all $n\geq a$ by {\em PMI}.''
\end{quote}
As long as someone can look back and see that you included the two necessary parts of the
proof, you do not necessarily need to point them out again. 
\end{enumerate}

\newpage
\section{Recursion}\label{sec:recursion}

You have seen examples of recursion if you have seen Russian Matryoshka dolls (Google it),
two almost parallel mirrors, a video camera pointed at the monitor, or a picture of a
painter painting a picture of a painter painting a picture of a painter...
More importantly for us, recursion is a very useful tool to implement algorithms.
You probably already learned about recursion in a previous programming course, but we 
present the concept in this brief section for the sake of review, and because it ties in
nicely with the other two topics in this chapter.\index{recursion}

\begin{df}\label{df:rec}
An algorithm is \hyperref[df:rec]{\em recursive} if it calls itself.\index{recursive}\index{recursion}
\end{df}

Examples of recursion that you may have already seen include {\em binary search}, 
{\em Quicksort}, and {\em merge sort}.

\begin{ques}\label{ques:ferzlerec}
Is following algorithm recursive?
Briefly explain.
\begin{code}
	int ferzle(int n) {
	    if(n<=0) {
	        return 3;
	    } else {
	        return ferzle(n-1) + 2;
	    }
	}
\end{code}
\answerlinetwo
\begin{answer}
Yes.  It clearly calls itself in the else clause.
\end{answer}
\end{ques}

If a subroutine/function simply called itself as a part of its execution, 
it would result in infinite recursion.  This is a bad thing.
Therefore, when using recursion, one must ensure that at some point,
the subroutine/function terminates without calling itself.
We will return to this point after we see what is perhaps 
the quintessential example of recursion.

\begin{exa}\label{exa:factRec}\index{factorial}
Notice that 
$$
\begin{array}{lllll}
0!&=&1 \\
1!&=&1 &=& 1 \times 0!\\
2!&=&2 \times 1 &=& 2 \times 1!\\
3!&=&3 \times 2 \times 1 &=& 3\times 2! \\
4!&=&4 \times 3\times 2\times 1 &=& 4\times 3!\\ 
5!&=&5\times 4 \times 3\times 2\times 1 &=& 5\times 4!\\ 
\multicolumn{5}{c}{\mbox{and in general, when } n>1,}\\
n!&=&n \times (n-1)\times\cdots\times 2\times 1 &=& n\times (n-1)!\\ 
\end{array}
$$

In other words, we can define $n!$ recursively as follows:

$$
n! = \left\{ \begin{array}{ll}
1 & \mbox{when $n = 0$}\\
 n * (n-1)! & \mbox{otherwise.}\\
\end{array}
\right.
$$

This leads to the following recursive algorithm to compute $n!$ when $n\geq 0$.

\begin{code}
    int factorial(int n) {
        if(n<=0) {
            return 1;
        } else {
            return n*factorial(n-1);
        }
    }
\end{code} 
\end{exa}

To guarantee that they will terminate, every recursive algorithm needs all
of the following.
\begin{enumerate}
\item
\index{base case (recursion)}
{\em Base case(s)}: One or more cases which are solved non-recursively.  
In other words, when an algorithm gets to the base case, it does not call itself again.
This is also called a {\em stopping case} or {\em terminating condition}.
\item 
{\em Inductive case(s)}: One or more recursive rule for all cases except the base case.
\index{inductive case (recursion)}
\item
{\em Progress:} The inductive case(s) should always progress toward the base case.
Often this means the arguments will get smaller until they approach the base case,
but sometimes it is more complicated than this.
\end{enumerate}
  
\vspace*{-.1in}

\begin{exa}
Let's take a closer look at the \scode{factorial} algorithm from Example~\ref{exa:factRec}.
If $n\leq 0$, \scode{factorial} does not make a recursive call.  Thus,
it has a {\em base case}.
When $n>0$, it is clearly making a recursive call, so it has {\em inductive cases}.
When a recursive call is made to \scode{factorial}, the
argument is smaller, so it is approaching a base case (i.e. making {\em progress}).
\end{exa}

\vspace*{-.1in}

\begin{ques}
Consider the \verb+ferzle+ algorithm from Question~\ref{ques:ferzlerec} above.
\begin{lenum}
\item What is/are the base case/cases?

\noindent
\answerline
\item What are the inductive cases?

\noindent
\answerline
\item Do the inductive cases make progress?

\noindent
\answerline
\end{lenum}
\begin{answer}
(a) The base cases are $n<\leq 0$.
(b) The inductive cases are $n>0$.
(c) Yes.  For any value $n>0$, the recursive call uses the value $n-1$, which is getting
closer to the base case of $0$.
\end{answer}
\end{ques}

\vspace*{-.1in}

\begin{exa}
Prove that the recursive \scode{factorial(n)} algorithm from Example~\ref{exa:factRec} 
returns $n!$ for all $n\geq 0$.
\begin{pf}
Notice that if $n=0$, \scode{factorial(0)} returns $1=0!$, so it works in that case.
For $k\geq 0$, assume \scode{factorial(k)} works correctly.
That is, it returns $k!$.
\scode{factorial(k+1)} return $k+1$ times the value returned by \scode{factorial(k)}.  
By the inductive hypothesis, \scode{factorial(k)} returns $k!$, so
\scode{factorial(k+1)} returns $(k+1)\times k!=(k+1)!$, as it should.
By PMI,  \scode{factorial(n)} returns $n!$ for all $n\geq 0$.
\end{pf}
\end{exa}

\begin{exa}\label{exa:countdown}
Implement an algorithm \scode{countdown(int n)} that outputs the integers from $n$ down to $1$, where $n>0$.
So, for example, \scode{countdown(5)} would output ``5 4 3 2 1''.
\begin{sol}
One way to do this is with a simple 
loop:
{\small
\begin{code}
    void countdown(int n) {
        for(i=n;i>0;i--) 
      	    print(i);	
    }
\end{code}
}
We wouldn't learn anything about recursion if we used this solution.
So let's consider how to do it with recursion.
Notice that \scode{countdown(n)} outputs $n$ followed by the numbers from $n-1$ down to $1$.
But the numbers $n-1$ down to $1$ are the output from \scode{countdown(n-1)}.
This leads to the following recursive algorithm:
{\small
\begin{code}
    void countdown(int n) {
        print(n);
        countdown(n-1):
    }
\end{code} 
}
To see if this is correct, we can trace through the
execution of \scode{countdown(3)}. The following table give 
the result.

\begin{center}
\begin{tabular}{|lcl|}
\hline
Execution of &outputs&then executes\\
\hline
\scode{countdown(3)}&3&\scode{countdown(2)}\\
\scode{countdown(2)}&2&\scode{countdown(1)}\\
\scode{countdown(1)}&1&\scode{countdown(0)}\\
\scode{countdown(0)}&0&\scode{countdown(-1)}\\
\scode{countdown(-1)}&-1&\scode{countdown(-2)}\\
\multicolumn{1}{|c}{$\vdots$}&$\vdots$&\multicolumn{1}{c|}{$\vdots$}\\
\hline
\end{tabular}
\end{center}
 
Unfortunately, \scode{countdown} will never terminate.  
We are supposed to stop printing when $n=1$, 
but we didn't take that into account.
In other words, we don't have a base case in our algorithm.
To fix this, we can modify it so that 
a call to \scode{countdown(0)} produces no output and does not call 
\scode{countdown} again.


Calls to \scode{countdown(n)} should also produce no output when $n<0$.
The following algorithm takes care of both problems and is our final solution.
{\small\begin{code}
    void countdown(int n) {
        if(n>0) {
            print(n);
            countdown(n-1);
        }
    }
\end{code}
}
Notice that when $n\leq 0$, \scode{countdown(n)} does nothing, 
making $n\leq 0$ the {\em base cases}.  
When $n>0$, \scode{countdown(n)} calls \scode{countdown(n-1)}, 
making $n>0$ the {\em inductive cases}.
Finally, when \scode{countdown(n)} makes a recursive call it is to \scode{countdown(n-1)}, 
so the inductive cases {\em progress} to the base case.
\end{sol}
\end{exa}

\newpage
\begin{exer}
Prove that the recursive \scode{countdown(n)} algorithm from Example~\ref{exa:countdown} 
works correctly.
(Hint:  Use induction.)

\vspace*{1.75in}
\begin{answer}
Notice that if $n\leq 0$, \scode{countdown(0)} prints nothing, so it works in that case.
For $k\geq 0$, assume \scode{countdown(k)} works correctly.\footnote{We are letting $n=0$ be the base case.  
You could also let $n=1$ be the base case, but then you would need to prove that \scode{countdown(1)} works.}  
Then \scode{countdown(k+1)} will print `$k+1$' 
and call \scode{countdown(k)}.  By the inductive hypothesis, 
\scode{countdown(k)} will print `k\, k-1\, $\ldots$ 2\, 1', so \scode{countdown(k+1)} will print `k+1\, k\, k-1 $\ldots$ 2\, 1',
so it works properly.  By PMI,  \scode{countdown(n)} works for all $n\geq 0$.
\end{answer}
\end{exer}


In general, we can solve a problem with recursion if we can: 
\begin{enumerate}
\item
Find one or more simple cases of the problem that can be solved directly.
\item
Find a way to break up the problem into smaller instances of the {\em same} problem.
\item
Find a way to combine the smaller solutions.
\end{enumerate}

Let's see a few classic examples of the use of recursion.

\begin{exa}\label{exa:binSearchComp}
Consider the {\em binary search} algorithm to find an item $v$ on a sorted list of size $n$.
The algorithm works as follows. 
\begin{itemize}
\item We compare the middle value $m$ of the array to $v$.
\item If the $m=v$, we are done.
\item Else if $m<v$, we binary search the left half of the array.
\item Else ($m>v$), we binary search the right half of the array.
\item Now, we have the same problem, but only half the size.
\end{itemize}
In Example~\ref{exa:binSearchIter} we saw the following iterative implementation of
binary search:
\begin{code}
    int binarySearch(int a[], int n, int val) {
       int left=0, right=n-1;
       while (right>=left) {
           int middle = (left+right)/2;
           if(val==a[middle]) 
               return middle;
           else if(val<a[middle]) 
               right=middle-1;
           else 
               left=middle+1; 
        }
       return -1;
    }
\end{code} 
Here is a version that uses recursion.  In this version we need to pass the 
endpoints of the array so we know what part of the array we are currently looking at.
\begin{code}
    int binarySearch(int[] a, int left, int right, int val) {
        if(right>=left) {
            int middle = (left+right)/2;
            if(val==a[middle])      
                return middle;
            else if(val<a[middle])  
                return binarySearch(a,left,middle-1,val);
            else                   
                return binarySearch(a,middle+1,right,val);
        } else {
            return -1;
        }
    }
\end{code}\index{binary search}\index{search!binary}
You should notice that in this case, the iterative and recursive algorithms are
very similar, and it is not clear that one implementation is better than the other.
However, if you were asked to write the algorithm from scratch, it is probably easier to
get the details right for the recursive one.
\end{exa}

\begin{exa}
Prove that the {\em recursive} \verb+binarySearch+ algorithm from
Example~\ref{exa:binSearchComp} is correct.
\begin{pf}
We will prove it by induction on $n=right-left+1$ (that is, the size of the array).

\noindent
{\bf Base case:}
If $n=0$, that means $right<left$, and \verb+binarySearch+ returns $-1$ as it should (since $val$ cannot possible be in an empty array).  
So it works
correctly for $n=0$.

\noindent
{\bf Inductive Hypothesis:}
Assume that \verb+binarySearch+ works for arrays of size $0$ through $k-1$ (we need strong induction for this proof).

\noindent
{\bf Inductive step:}
Assume \verb+binarySearch+ is called on an array of size $k$.  There are three cases.
\begin{itemize}
\item
If $val=a[middle]$, the algorithm returns $middle$ which is the correct answer.
\item
If $val<a[middle]$, a recursive call is made on the first half of the array (from $left$ to $middle-1$).  
Because $a$ is sorted, if $val$ is in the array, it is in that half of the array, so we just need to prove that the
recursive call returns the correct value.
Notice that the first half of the array has less than $n$ elements (it does not contain $middle$ or anything to the right of $middle$, so
it is clearly smaller by at least one element).
Thus, by the inductive hypothesis, it returns the correct index or $-1$ if $val$ is not in that part of the array.  
Therefore it returns the correct value.
\item
The case for $val>a[middle]$ is symmetric to the previous case and the details are left to the reader.
\end{itemize}
In all cases, it works correctly on an array of size $k$.

\noindent
{\bf Summary:}
Since it works for an array of size $0$ and whenever it works for arrays of size at most $k-1$ it works for arrays of size $k$,
by the principle of mathematical induction, it works for arrays of any nonnegative size.
\end{pf}
\end{exa}

\begin{rem}
You might think the base case in the previous proof should be $n=1$, but that is not actually correct. A failed search will always
make a final call to \verb+binarySearch+ with $n=0$.  If we don't prove it works for an empty array then we cannot
be certain that it works for failed searches.
\end{rem}

\begin{exa}\label{exa:fibAlg}\index{Fibonacci numbers}
Recall the {\em Fibonacci sequence}, defined by the recurrence relation
$$f_n=\left\{
\begin{array}{ll}
0&\mbox{ if $n$=0}\\
1&\mbox{ if $n$=1}\\
f_{n-1}+f_{n-2}&\mbox{ if }n>1.\\
\end{array}
\right.
$$
Let's see an iterative and a recursive algorithm to compute $f_n$.
The iterative algorithm (on the left) starts with $f_0$ and $f_1$ and computes
each $f_i$ based on $f_{i-1}$ and $f_{i-2}$ for $i$ from $2$ to $n$.
As it goes, it needs to keep track of the previous two values.
The recursive algorithm (on the right) just uses the definition and is pretty
straightforward.

\noindent
\begin{minipage}[t]{.4\textwidth}
\begin{code}
int Fib(int n) {
    int fib;
    if(n <= 1) {
       return(n);
    } else {
       int fibm2=0;
       int fibm1=1;
       int index=1;
       while(index < n) {
          fib=fibm1+fibm2;
          fibm2=fibm1;
          fibm1=fib;
          index++;
       }
       return(fib); 
    }
}
\end{code} 
\end{minipage}
\hfill 
\begin{minipage}[t]{.5\textwidth}
\begin{code}
int FibR(int n) {
    if(n <= 1) {
       return(n);
    } else { 
       return(FibR(n-1)+FibR(n-2));
    }
}
\end{code} 
\end{minipage}
\end{exa}

\begin{ques}\label{ques:fibrineff}
Which algorithm is better, \verb+Fib+ or \verb+FibR+? Give several reasons to justify your answer.

\noindent
\answerlinethree
\begin{answer}
It is pretty clear that the recursive algorithm
is much shorter and was a lot easier to write.  
It is also a lot easier to
make a mistake implementing the iterative algorithm.
So far, it looks like the recursive algorithm is the clear winner.
However, in the next section we will show you why the recursive 
algorithm we gave should never be implemented.  It turns out that is is
very inefficient.

The bottom line is that the iterative algorithm is better in this case.
Don't feel bad if you thought the recursive algorithm was better.
After the next section, you will be better prepared to compare
recursive and iterative algorithms in terms of efficiency.
\end{answer}
\end{ques}

Although recursion is a great technique to solve many problems, care must be taken when 
using it.  It easy to make simple mistakes like we did in Example~\ref{exa:countdown}.
They can also be very inefficient on occasion, as we alluded to in the previous example
(and will prove later).
In addition, recursive algorithms often take more memory than iterative ones, as we will see next.

\begin{exa}
Consider our algorithms for $n!$.  
The iterative one from Example \ref{exa:factorial} uses memory to store four numbers: 
$n$, $f$, $i$, and return value.\footnote{I won't get technical here,
but memory needs to be allocated for the value returned by a function.}
The recursive one from Example \ref{exa:factRec} uses memory to store two numbers: $n$ and the
return value.
Although the recursive algorithm uses less memory, it is called multiple times, and every 
call needs its own memory.  
For instance, a call to \scode{factorial(3)} will
call \scode{factorial(2)} which will call \scode{factorial(1)}.  
Thus, computing $3!$ requires enough memory to store $6$ numbers, which is more than the $4$ required by the
iterative algorithm.  In general, the recursive algorithm to compute $n!$
will need to store $2n$ numbers, whereas the iterative one will still just need $4$, no matter
how large $n$ gets.
\end{exa}

Since computers have a finite amount of memory, and since every call to a function requires its own memory,
there is a limit to how many recursive calls can be made in practice.  
In fact some languages, including Java, have a defined limit of how deep the recursion can be.
Even for those that don't have a limit, if you run out of memory, you can certainly expect bad things to happen.
This is one of the reasons recursion is avoided when possible. 

Good compilers attempt to remove recursion, but it is not always possible.
Good programmers do the same.  Since recursive algorithms are often more intuitive, it often makes sense
to think in terms of them.  But many recursive algorithms can be turned into iterative algorithms
that are as efficient and use less memory.  There is no single technique to
do so, and it is not always necessary, but it is a good thing to keep in mind.

Let's see a few more examples of the subtle problems that we can run into when using recursion.

\begin{exa}
The following algorithm is supposed to sum the numbers from $1$ to $n$:
{\small\begin{code}
    void Sum1toN(int n) {
        if (n == 0)  return(0);
        else         return(n + Sum1toN(n-1));
    }
\end{code} 
}
Although this algorithm works fine for non-negative values of $n$, 
it will go into infinite recursion if $n<0$.
Like our original solution to the \scode{countdown} problem, the
mistake here is an {\em improper base case}.
\end{exa}

It is easy to get things backwards when recursion is involved.
Consider the following example.

\begin{ques}
One of these routines prints from 1 up to $n$,
the other from $n$ down to 1.
Which does which?

\begin{minipage}{.4\textwidth}
\begin{code}
     void PrintN(int n) { 
          if (n > 0) {
             PrintN(n-1);
             print(n);
          }                          
     }
\end{code}
\end{minipage}
\hfill 
\begin{minipage}{.4\textwidth}
\begin{code}
     void NPrint(int n) {
         if (n > 0) {
            print(n);
            NPrint(n-1);
         }
     }
\end{code} 
\end{minipage}

\answerlinetwo
\begin{answer}
\verb+PrintN+ will print from $1$ to $n$, and \verb+NPrint+ will print from $n$ to $1$.
If you go the answer wrong, go back and convince yourself that this is correct.
\end{answer}
\end{ques}


We conclude this section by summarizing some of the advantages and disadvantages
of recursion.

The advantages include:
\begin{enumerate}
\item 
Recursion often mimics the way we think about a problem, thus the
recursive solutions can be very intuitive to program.
\item 
Often recursive algorithms to solve problems are much shorter than
iterative ones. This can  make the code easier to understand, modify, and/or debug.
\item  
The best known algorithms for many problems are based on
a divide-and-conquer approach:
\begin{itemize}
\item  Divide the problem into a set of smaller problems
\item  Solve each small problem separately
\item  Put the results back together for the overall solution
\end{itemize}
These divide-and-conquer techniques are often best thought of in terms of
recursive algorithms. 
\end{enumerate}

Perhaps the main disadvantage of recursion is the extra time and space
required.
We have already discussed the extra space.  
The extra time comes from the fact that when a recursive call is made,
the operating system has to record how to restart the calling subroutine later on,
pass the parameters from the calling subroutine to the called
subroutine (often by pushing the parameters onto a stack controlled by the system),
set up space for the called subroutine's local variables, etc.
The bottom line is that calling a function is not ``free''.  

Another disadvantage is the fact that sometimes a slick-looking recursive algorithm 
turns out to be very inefficient.  We alluded to this in Example~\ref{ques:fibrineff}.
On the other hand, if such inefficiencies are found, there are techniques
that can often easily remove them 
(e.g. a technique called memoization\footnote{No, that's not a typo.  Google it.}).
But you first have to remember to analyze 
your algorithm to determine whether or not there might be an efficiency problem.

%%%%%%%%%%%%%%%%%%%%%%%%%%%%%%%%%%%%%%%%%%%%%%%%%%%%%%%%%%%%%%%%%%%%%%%%%%%%
%%%%%%%%%%%%%%%%%%%%%%%%%%%%%%%%%%%%%%%%%%%%%%%%%%%%%%%%%%%%%%%%%%%%%%%%%%%%
%%%%%%%%%%%%%%%%%%%%%%%%%%%%%%%%%%%%%%%%%%%%%%%%%%%%%%%%%%%%%%%%%%%%%%%%%%%%
%%%%%%%%%%%%%%%%%%%%%%%%%%%%%%%%%%%%%%%%%%%%%%%%%%%%%%%%%%%%%%%%%%%%%%%%%%%%
%%%%%%%%%%%%%%%%%%%%%%%%%%%%%%%%%%%%%%%%%%%%%%%%%%%%%%%%%%%%%%%%%%%%%%%%%%%%
\newpage

\section{Solving Recurrence Relations}\label{sec:recRel}
Recall that a {\em recurrence relation}\index{recurrence relations}
is simply a sequence that is recursively defined.
More formally, a recurrence relation is a formula that defines $a_n$ 
in terms of $a_i$, for one or more values of $i<n$.\footnote{You might also see recurrence relations
written using function notation, like $a(n)$.
Although there are technical differences between these notations, you can think of them as
being essentially equivalent in this context.}

\begin{exa}\index{Fibonacci numbers}
We previously saw that we can define $n!$ by $0!=1$, and if $n>0$, $n!=n\cdot(n-1)!$.
This is a recurrence relation for the sequence $n!$.

Similarly, we have seen the Fibonacci sequence several times.
Recall that $n$-th Fibonacci number is given by $f_0=f_1=1$ and for $n>1$, 
$f_n=f_{n-1}+f_{n-2}$.  This is recurrence relation for the sequence of Fibonacci numbers.
\end{exa}

\vspace*{-.05in} % For later exercise that wants to go to next page.

\begin{exa}\label{exa:somerecs}
Each of the following are recurrence relations.
\begin{eqnarray*}
t_n &=& n\cdot t_{n-1} + 4\cdot t_{n-3}\\
r_n &=& r_{n/2}+1\\
a_n &=& a_{n-1}+2\cdot a_{n-2}+3\cdot a_{n-3}+4\cdot a_{n-4}\\
p_n &=& p_{n-1}\cdot p_{n-2}\\
s_n &=& s_{n-3} + n^2 - 4 n + 32
\end{eqnarray*}
We have not given any initial conditions for these recurrence relations.
Without initial conditions, we cannot compute particular values.
We also cannot solve the recurrence relation uniquely.
\end{exa}

Recurrence relations have 2 types of terms: {\em recursive} term(s) and
the {\em non-recursive} terms.
\index{non-recursive term (recurrence relation)}\index{recursive term (recurrence relation)}
%% Huh?  This isn't correct--the base cases are independent of the recursive definition.
%% Where did I get this idea?  CAC, 6/11/14.
%These are synonymous with the {\em inductive} and {\em base} cases of recursive algorithms.
In the previous example, the recursive term of $s_n$ is $s_{n-3}$
and the non-recursive term is  $n^2 - 4 n + 32$.  

\begin{ques}
Consider the recurrence relations $r_n$ and $a_n$ from Example~\ref{exa:somerecs}.
\begin{lenum}
\item
What are the {\em recursive} terms of $r_n$?

\noindent
\answerline
\item
What are the {\em non-recursive} terms of $r_n$?

\noindent
\answerline
\item
What are the {\em recursive} terms of $a_n$?

\noindent
\answerline
\item
What are the {\em non-recursive} terms of $a_n$?

\noindent
\answerline
\end{lenum}
\begin{answer}
(a) $r_{n/2}$.
(b) $1$.
(c) $a_{n-1}+2\cdot a_{n-2}+3\cdot a_{n-3}+4\cdot a_{n-4}$.
(d) There are none.
\end{answer}
\end{ques}

In computer science, the most common place we use recurrence relations is to analyze recursive algorithms.
We won't get too technical yet, but let's see a simple example.

\begin{exa}\label{exa:factMuls}
How many multiplications are required to compute $n!$ using the \verb+factorial+ algorithm given
in Example \ref{exa:factRec} (repeated below)?
\begin{code}
    int factorial(int n) {
        if(n<=0) {
            return 1;
        } else {
            return n*factorial(n-1);
        }
    }
\end{code} 
\begin{sol}
Let $M_n$ be the number of multiplications needed to compute $n!$ using 
the \verb+factorial+ algorithm from Example \ref{exa:factRec}.
From the code, it is obvious that $M_0=0$.  If $n>0$, 
the algorithm uses one multiplication and then makes a recursive call to \verb+factorial(n-1)+.
By the way we defined $M_n$, \verb+factorial(n-1)+ does $M_{n-1}$ multiplications.
Therefore, $M_n=M_{n-1}+1$.  

So the recurrence relation for the number of multiplications is 
$$M_n=\left\{
\begin{array}{ll}
0&\mbox{ if $n$=0}\\
M_{n-1}+1&\mbox{ if }n>0.\\
\end{array}
\right.
$$
\end{sol}
\end{exa}

Given a recurrence relation for $a_n$, you can't just plug in $n$ and get an answer.  For instance, 
if $a_n=n\cdot a_{n-1}$, and $a_1=1$, what is $a_{100}$?  The only obvious way to compute it is
to compute $a_2, a_3,\ldots, a_{99}$, and then finally $a_{100}$.  That is the reason why
{\em solving} recurrence relations is so important.  As mentioned previously, solving a recurrence
relation simply means finding a {\em closed form expression} for it.
\index{closed form (recurrence relation)}
\index{recurrence relations!solving}

\begin{exa}
It is not too difficult to see that the recurrence from Example \ref{exa:factMuls} has the
solution $M_n=n$.  To prove it, notice that with this assumption, $M_{n-1}+1= (n-1)+1 = n=M_{n}$,
so the solution is consistent with the recurrence relation.

We can also prove it with induction: 
We know that $M_0=0$, so the base case of $k=0$ is true.
Assume $M_k=k$ for $k\ge 0$.  Then we have
$$M_{k+1} = M_{k} + 1 = k + 1,$$
so the formula is correct for $k+1$.  Thus, by PMI, the formula is correct for all $k\geq 0$.
\end{exa}

The last example demonstrates an important fact about recurrence relations used to analyze algorithms.
The recursive terms come from when a recursive function calls itself.  
The non-recursive terms come from the other work that is done by the function,
including any splitting or combining of data that must be done.

\newpage

\begin{exa}\label{exa:binSearchRecRel}
\index{binary search}\index{search!binary}
Consider the recursive {\em binary search} algorithm we saw in Example~\ref{exa:binSearchComp}:
\begin{code}
    int binarySearch(int[] a, int left, int right, int val) {
        if(right>=left) {
            int middle = (left+right)/2;
            if(val==a[middle])      
                return middle;
            else if(val<a[middle])  
                return binarySearch(a,left,middle-1,val);
            else                   
                return binarySearch(a,middle+1,right,val);
        } else {
            return -1;
        }
    }
\end{code}
Find a recurrence relation for the worst-case complexity of \scode{binarySearch}.
\begin{sol}
Let $T_n$ be the complexity of \scode{binarySearch} for an array of size $n$.
Notice that the only things done in the algorithm are to find the middle element,
make a few comparisons, perhaps make a recursive call, and return a value.  Aside from
the recursive call, the amount of work done is constant, which we will just call $1$ operation.  
Notice that at most one recursive call is made, and that the array passed in is half
the size.
Therefore  $T_n=T_{n/2}+1$.\footnote{Technically, 
the recurrence relation is  $T_n=T_{\lfloor n/2\rfloor}+1$ since $n/2$ might not be an integer. 
It turns out that most of the time we can ignore the floors/ceilings and still obtain the correct answer.
}
If we want a base case, we can use $T_0=1$ since the algorithm will simply return $-1$ for an
empty array, and that clearly takes constant time.
We'll see how to solve this recurrence shortly.
\end{sol}
\end{exa}

We will discuss using recurrence relations to analyze recursive algorithms in more detail
in section~\ref{sec:anaRecAlg}.  But first we will discuss how to solve recurrence relations.
There is no general method to solve recurrences.
There are many strategies, however.  In the next few sections we will discuss
four common techniques: 
the {\em substitution method},  
the {\em iteration method},
the {\em Master Theorem} (or {\em Master Method}),
and the {\em characteristic equation method} for linear recurrences.

\begin{ques}
Let's see if you have been paying attention.  
What does it mean to solve a recurrence relation?

\noindent
\answerlinetwo
\begin{answer}
It means to find a closed-form expression for it.  In other words, one that does not
define the sequence recursively.
\end{answer}
\end{ques}

As we continue our discussion of recurrence relations, you will notice that we will begin to
sometimes use the function notation (e.g. $T(n)$ instead of $T_n$).  
We do this for several reasons.  The first is so that you are comfortable with either notation.
The second is that in algorithm analysis, this notation seems to be more common,
at least in my experience.

\subsection{Substitution Method}
The {\em substitution method} might be better called the 
{\em guess and prove it by induction method}. 
\index{substitution method}
\index{recurrence relations!solving!substitution method} 
Why?  Because to use it, you first have to figure out what you think the
solution is, and then you need to actually prove it.  Because of the close
tie between recurrence relations and induction, it is the most natural technique to use.
Let's see an example.
\begin{exa}
Consider the recurrence 
$$
S(n) = \left\{ 
\begin{array}{ll}
1&\mbox{when $n = 1$}\\
S(n-1) + n & \mbox{otherwise}
\end{array}
\right.
$$
Prove that the solution is 
${\displaystyle S(n)=\frac{n(n+1)}{2}.}$

\medskip
\noindent
\begin{pf}
When $n=1$, $S(1)=1 = \frac{1(1+1)}{2}$.
Assume that $S(k-1)=\frac{(k-1)k}{2}$.
Then
\begin{eqnarray*}
S(k)&=&S(k-1)+k\, \, \, \mbox{(Definition of }S(k)\mbox{)}\\
&=&\dfrac{(k-1)(k)}{2}+k\, \, \, \mbox{(Inductive hypothesis)}\\
&=&\dfrac{k^2-k}{2}+k\, \, \, \mbox{(The rest is just algebra)}\\
&=&\dfrac{k^2-k+2 k}{2}\\
&=&\dfrac{k^2+k}{2}\\
&=&\dfrac{k(k+1)}{2}.
\end{eqnarray*}
By PMI, $S(n)=\frac{n(n+1)}{2}$ for all $n\geq 1$.
\end{pf}
\end{exa}

\begin{exer}
Recall that in Example~\ref{exa:binSearchRecRel}, we
developed the recurrence relation $T(n)=T(n/2)+1, T(0)=1$ for the complexity of \verb+binarySearch+.
For technical reasons, ignore $T(0)$ and assume $T(1)=1$ is the base case.
Use substitution to prove that $T(n)=\log_2 n + 1$ is a solution to this 
recurrence relation.

\vspace*{2.5in}
\begin{answer}
When $n=1$, $T(1)=1=0+1=\log_2 1 + 1$.
Assume that $T(k)=\log_2 k + 1$ for all $1\leq k <n$ 
(we are using strong induction).
Then 
\begin{eqnarray*}
T(n)&=&T(n/2)+1\\
&=&(\log_2 (n/2) + 1) + 1\\
&=&\log_2 n - \log_2 2 + 2\\
&=&\log_2 n - 1 + 2\\
&=&\log_2 n +1.
\end{eqnarray*}
So by PMI, $T(n)=\log_2 n + 1$ for all $n\geq 1$.
\end{answer}
\end{exer}

\begin{exa}\label{exa:Hanoi1}
Solve the recurrence
$$
H_n = \left\{ 
\begin{array}{ll}
1&\mbox{when $n = 1$}\\
2 H_{n-1} + 1 & \mbox{otherwise}
\end{array}
\right.
$$
\begin{pf}
Notice that $H_1=1$, $H_2=2\cdot 1 + 1 = 3$, $H_3=2\cdot 3+1 = 7$, and $H_4=2\cdot 7 + 1=15$.
It sure looks like $H_n=2^n-1$,  but now we need to prove it.  
Since $H_1=1=2^1-1$, we have our base case of $n=1$.  Assume $H_n=2^n-1$.  Then
\begin{eqnarray*}
H_{n+1}&=&2H_n+1\\
&=&2(2^n-1)+1\\
&=&2^{n+1}-1,
\end{eqnarray*}
and the result follows by induction.
\end{pf}
\end{exa}

\begin{exer}
Solve the following recurrence relation and use induction to prove
your solution is correct:  $A(n)=A(n-1)+2$, $A(1)=2$.

\vspace*{4.5in}
\begin{answer}
We begin by computing a few values to see if we can find a pattern.
$A(2)=A(1)+2 = 2+2= 4$, $A(3)=A(2)+2=4+2=6$, $A(4)=8$, $A(5)=10$, etc.
It seems pretty obvious that $A(n)=2n$.  It holds for $n=1$, so we have our base case.
Assume $A(n)=2n$.  Then $A(n+1)=A(n)+2 = 2n+2=2(n+1)$, so it holds for $n+1$.
By PMI, $A(n)=2n$ for all $n\geq 1$.
\end{answer}
\end{exer}

\newpage
\begin{exa}\index{Fibonacci numbers}
Why was the recursive algorithm to compute $f_n$ from Example \ref{exa:fibAlg} so bad?
\begin{sol}
Let's count the number of additions \scode{FibR(n)} computes since that is the main
thing that the algorithm does.\footnote{Alternatively, we could count the number of recursive calls made.
This is reasonable since the amount of work done by the algorithm, aside from the recursive calls,
is constant.  Therefore, the time it takes to compute $f_n$ is proportional to 
the number of recursive calls made.
This would produce a slightly different answer, but they would be comparable.}
%Footnote causes an undefull vbox, but I don't see how to get rid of it.
Let $F(n)$ be the number of additions required to compute $f_n$ using \scode{FibR(n)}.
Since \scode{FibR(n)} calls \scode{FibR(n-1)} and \scode{FibR(n-2)}
and then performs one addition,
it is easy to see that 
$$F(n)=F(n-1)+F(n-2)+1,$$
where $F(0)=F(1)=0$ is clear from the algorithm.
We could use the method for linear recurrences that will be outlined later to solve this, 
but the algebra gets a bit messy.  
Instead, Let's see if we can figure it out by computing some values.
\begin{eqnarray*}
F(0)&=&0\\
F(1)&=&0\\
F(2)&=&F(1)+F(0)+1 = 1\\
F(3)&=&F(2)+F(1)+1 = 2\\
F(4)&=&F(3)+F(2)+1 = 4\\
F(5)&=&F(4)+F(3)+1 = 7\\
F(6)&=&F(5)+F(4)+1 = 12\\
F(7)&=&F(6)+F(5)+1 = 20
\end{eqnarray*}

No pattern is evident unless you add one to each of these. If you do, you will
get $1,1,2,3,5,8,13,21$, etc., which looks a lot like the Fibonacci sequence starting with
$f_1$.  So it appears $F(n)=f_{n+1}-1$.
To verify this, first notice that $F(0)=0=f_1-1$ and $F(1)=0=f_2-1$.
Assume it holds for all values less than $k$.
Then
\begin{eqnarray*}
F(k)&=&F(k-1)+F(k-2)+1\\
&=&f_{k}-1+f_{k-1}-1 + 1\\
&=&f_{k}+f_{k-1} -1\\
&=&f_{k+1}-1.
\end{eqnarray*}
%\vspace*{0in} % gets rids of \vbox underfull warning.

\noindent
The result follows by induction.

So what does this mean?  It means in order to compute $f_n$, 
\scode{FibR(n)} performs $f_{n+1}-1$ additions.
In other words, it computes $f_n$ by adding a bunch of $0$s and $1$s,
which doesn't seem very efficient.  
Since $f_n$ grows exponentially (we'll see this in Example \ref{exa:fibSol}),
then $F(n)$ does as well.  That pretty much explains what is wrong with the recursive algorithm.
\end{sol}

\end{exa}

\subsection{Iteration Method} 

With the iteration method (sometimes called {\em backward substitution}, 
we expand the recurrence and express
it as a summation dependent only on $n$ and initial conditions.
Then we evaluate the summation.
Sometimes the closed form of the sum is obvious as we are iterating
(so no actual summation appears in our work), while at other times
it is not (in which case we {\em do} end up with an actual summation).
\index{iteration method}
\index{recurrence relations!solving!iteration method} 

Our first example perhaps has too many steps of algebra, but it never
hurts to be extra careful when doing so much algebra.
We also don't provide a whole lot of justification or explanation for
the steps.  We will do that in the next example.  It is easier to see
the overall idea of the iteration method if we don't interrupt it
with comments.  If this example does not make sense, come back to it
after reading the next example.

\begin{exa}\label{exa:R}
Solve the recurrence 
$$
R(n) = \left\{ 
\begin{array}{ll}
1&\mbox{when $n = 1$}\\
2 R(n/2) + n/2 & \mbox{otherwise}
\end{array}
\right.
$$
\begin{pf}
We have
\begin{eqnarray*}
R(n)&=&2 R(n/2) + n/2\\
&=&2(2 R(n/4) + n/4 ) + n/2\\
&&=2^2R(n/4) + n/2+n/2\\
&&= 2^2R(n/4) + n\\
&=&2^2(2 R(n/8) + n/8 ) + n\\
&&=  2^3R(n/8) + n/2+n\\
&&=  2^3R(n/8) + 3n/2\\
&=&2^3(2 R(n/16) + n/16 ) + 3n/2\\
&&=2^4R(n/16) + n/2 + 3n/2\\
&&=2^4R(n/16) + 2n\\
&\hfill\vdots\hfill\hfill\\
&=&2^k R(n/(2^k)) + k n/2\\
&=&2^{\log_2 n} R(n/(2^{\log_2 n})) + (\log_2 n) n/2\\
&=&n R(n/n) + (\log_2 n) n/2\\
&=&n R(1) + (\log_2 n) n/2\\
&=&n + (\log_2 n) n/2
\end{eqnarray*}
\end{pf}\end{exa}


Using this method requires a little abstract thinking and pattern recognition.
It also requires good algebra skills.  
Care must be taken when doing algebra, especially with the non-recursive terms. 
Sometimes you should add/multiply  (depending on context) them all together,
and other times you should leave them as is.  The problem is that it takes
experience (i.e. practice) to determine which one is better in a given situation.
The key is flexibility.  If you try doing it one way and don't see a pattern, 
try another way.

Here is my suggestion for using this method
\begin{enumerate}
\item
Iterate enough times so you are certain of what the pattern is.  Typically this means
at least 3 or 4 iterations.
\item
As you iterate, make adjustments to your algebra as necessary so you can see the pattern.
For instance, whether you write $2^3$ or $8$ can make a difference in seeing the pattern.
\item
Once you see the pattern, generalize it, writing what it should look like after $k$ iterations.
\item
Determine the value of $k$ that will get you to the base case, and then plug it in.
\item
Simplify.
\end{enumerate}

\begin{ques}
The {\em iteration method} is probably not a good choice to solve the following recurrence relation.
Explain why. 
$$T(n)=T(n-1)+3T(n-2)+n*T(n/3)+n^2,\, \, \,  T(1)=17$$
\answerlinethree
\begin{answer}
It contains 3 very different looking recursive terms so it is very unlikely we will
be able to find any sort of meaningful pattern by iteration.
\end{answer}
\end{ques}

Here is an example that contains more of an explanation of the technique.
\begin{exa}\label{exa:longiterone}
Solve the recurrence relation $T(n)=2T(n/2)+ n^3$, $T(1)=1$.
\begin{sol}
We start by backward substitution:
\begin{eqnarray*}
T(n)&=&2T(n/2)+n^3\\
    &=&2[2T(n/4)+(n/2)^3]+n^3\\
    &=&2[2T(n/4)+n^3/8)]+n^3\\
    &=&2^2T(n/4)+n^3/4+n^3
\end{eqnarray*}
Notice that in the second line we have $(n/2)^3$ and not $n^3$.  
This may be more clear if rewrite the formula using $k$: $T(k)=2T(k/2)+k^3$.
When applying the formula to $T(n/2)$, we have $k=n/2$, so we get
$$T(n/2)=2T((n/2)/2)+(n/2)^3=2T(n/4)+n^3/8.$$
Back to the second line, 
also notice that the $2$ is multiplied by both the $2T(n/4)$ and the $(n/2)^3$ terms.
A common error is to lose one of the $2$s on the $T(n/4)$ term or miss it on
the $(n/2)^3$ term when simplifying.
Also, $(n/2)^3=n^3/8$, not $n^3/2$.  This is another common mistake.
Continuing,

\begin{eqnarray*}
T(n)&=&\ldots\\
    &=&2^2T(n/4)+n^3/4+n^3\\
    &=&2^2[2T(n/8)+(n/4)^3]+n^3/4+n^3\\
    &=&2^2[2T(n/8)+n^3/4^3]+n^3/4+n^3\\
    &=&2^3T(n/8)+n^3/4^2+n^3/4+n^3.
\end{eqnarray*}
By now you should have noticed that I use 2 or more steps for every iteration--I do one substitution
and then simplify it before moving on to the next substitution.  This helps to 
ensure I don't make algebra mistakes and that I can write it out in a way that 
helps me see a pattern.

Next, notice that we can write the last line as
$$2^3T(n/2^3)+n^3/4^2+n^3/4^1+n^3/4^0,$$
so it appears that we can generalize this to 
$$2^kT(n/2^k)+\sum_{i=0}^{k-1} n^3/4^i.$$
The sum starts at $i=0$ (not $1$) and goes to $k-1$ (not $k$).  
It is easy to get either (or both) of these wrong if you aren't careful.
We should be careful to make sure we have seen the correct pattern.
Too often I have seen students make a pattern out of 2 iterations.
Not only is this not enough iterations to be sure of anything, 
the pattern they usually come up with only holds for the last iteration they did. 
The pattern has to match {\em every} iteration.  
To be safe, go one more iteration after you
identify the pattern to verify that it is correct.

Continuing (with a few more steps shown to make all of the algebra as clear as possible), we get
\begin{eqnarray*}
T(n)&=&\ldots\\
    &=&2^3T(n/2^3)+n^3/4^2+n^3/4^1+n^3/4^0\\
    &=&2^3[2T(n/2^4)+(n/2^3)^3]+n^3/4^2+n^3/4^1+n^3/4^0\\
    &=&2^3[2T(n/2^4)+n^3/2^9]+n^3/4^2+n^3/4^1+n^3/4^0\\
    &=&2^4T(n/2^4)+n^3/2^6+n^3/4^2+n^3/4^1+n^3/4^0\\
    &=&2^4T(n/2^4)+n^3/4^3+n^3/4^2+n^3/4^1+n^3/4^0\\
    &=&\ldots\\
    &=&2^kT(n/2^k)+\sum_{i=0}^{k-1} n^3/4^i.
\end{eqnarray*}
Notice that this {\em does} seem to match the pattern we saw above.
We can evaluate the sum to simplify it a little more:

\begin{eqnarray*}
T(n)&=&\ldots\\
    &=&2^kT(n/2^k)+\sum_{i=0}^{k-1} n^3/4^i\\
    &=&2^kT(n/2^k)+n^3\sum_{i=0}^{k-1} 1/4^i\\
    &=&2^kT(n/2^k)+n^3\sum_{i=0}^{k-1} (1/4)^i\\
    &=&2^kT(n/2^k)+n^3\left(\frac{1-(1/4)^k}{1-1/4}\right)\\
    &=&2^kT(n/2^k)+n^3(4/3)(1-(1/4)^k)\\
\end{eqnarray*}
We are almost done.  We just need to find a $k$ that allows us to get rid of the recursion.
Thus, we need to determine what value of $k$ makes $T(n/2^k)=T(1)=1$.  In other words, we need
$k$ such that  
$$n/2^k=1.$$
This is equivalent to 
$$n=2^k.$$
Taking log (base 2) of both sides, we obtain 
$$\log_2 n=\log_2(2^k)=k\log_{2} 2=k.$$  
So $k=\log_2 n$.
We plug in $k$ and use the fact that $2^{\log_2 n}=n$ along with the exponent rules to obtain
\begin{eqnarray*}
T(n)&=&\ldots\\
    &=&2^kT(n/2^k)+n^3(4/3)(1-(1/4)^k)\\
    &=&2^{\log_2 n}T(n/2^{\log_2 n})+n^3(4/3)(1-(1/4)^{\log_2 n})\\
    &=&nT(1)+n^3(4/3)\left(1-\frac{1}{(2^2)^{\log_2 n}}\right)\\
    &=&n\cdot 1+n^3(4/3)\left(1-\frac{1}{(2^{\log_2 n})^2}\right)\\
    &=&n+n^3(4/3)\left(1-\frac{1}{n^2}\right)\\
    &=&n+\frac{4}{3}n^3-\frac{4}{3}n\\
    &=&\frac{4}{3}n^3-\frac{1}{3}n.
\end{eqnarray*}
So we have that $T(n)=\frac{4}{3}n^3-\frac{1}{3}n$.
\end{sol}
\end{exa}

\begin{exer}\label{exa:Hanoi2}
Use iteration to solve the recurrence
$$
H(n) = \left\{ 
\begin{array}{ll}
1&\mbox{when $n = 1$}\\
2 H(n-1) + 1 & \mbox{otherwise}
\end{array}
\right.
$$
\vspace*{7.5in}
\begin{answer} 
\begin{eqnarray*}
H(n)&=&2 H(n-1) + 1\\
&=&2(2 H(n-2) + 1) + 1\\
&&= 2^2 H(n-2) + 2 + 1\\
&=&2^2(2 H(n-3) + 1) + 2 + 1\\
&&= 2^3 H(n-3) + 2^2 + 2 + 1\\
&\vdots&\\
&=&2^{n-1}H(1) + 2^{n-2} + 2^{n-3}+\cdots + 2 + 1\\
&=&2^{n-1} + 2^{n-2} + 2^{n-3}+\cdots + 2 + 1\\
&=&2^n-1
\end{eqnarray*}
Thus, $H(n)=2^n-1$. Luckily, this matches our answer from Example~\ref{exa:Hanoi1}. 
\end{answer}
\end{exer}

\newpage

\begin{exa}
Give a tight bound for the recurrence $T(n)=T(\sqrt{n})+1$, where $T(2)=1$.
\begin{sol}
We can see that 
\begin{eqnarray*}
T(n)&=&T(n^{1/2})+1\\
&=&T(n^{1/4})+1+1\\
&=&T(n^{1/8})+1+1+1\\
&=&T(n^{1/2^k})+k\\
\end{eqnarray*}
If we can determine when $n^{1/2^k}=2$, we can obtain a solution.
Taking logs (base 2) on both sides, we get 
$$\log_2(n^{1/2^k})=\log_2 2.$$ 
We apply the power-inside-a-log rule
and the fact that $\log_2 2=1$ to get
$$(1/2^k)\log_2 n=1.$$
Multiplying both sides by $2^k$ and flipping it around, we get
$$2^k= \log_2 n.$$
Again taking logs, we get 
$$k=\log_2\log_2 n.$$
Therefore, 
\begin{eqnarray*}
T(n)&=&T(n^{1/2^{\log_2\log_2 n}})+\log_2\log_2 n\\
     &=&T(2)+\log_2\log_2 n\, \, \, \mbox{(since $n^{1/2^{\log_2\log_2 n}}=2$ by the way we chose $k$)}\\
     &=&1+\log_2\log_2 n.\\
\end{eqnarray*}
Therefore, $T(n)=1+\log_2\log_2 n$.
\end{sol}
\end{exa}

\newpage
\begin{exer}
Use iteration to solve the recurrence relation that we developed in Example~\ref{exa:binSearchRecRel}
for the complexity of \verb+binarySearch+:
$$T(n)=T(n/2)+1, T(1)=1.$$
\vspace*{7.5in}
\begin{answer}
Iterating a few steps, we discover:
\begin{eqnarray*}
T(n)&=&T(n/2)+1\\
&=&T(n/4)+1+1\\
&=&T(n/2^2)+2\, \, \, \, \mbox{(I think I see a pattern!)}\\
&=&T(n/2^3)+1+2\\
&=&T(n/2^3)+3\, \, \, \, \mbox{(I {\em do} see a pattern!)}\\
&\vdots&\\
&=&T(n/2^k)+k\\
\end{eqnarray*}
We need to find $k$ such that $n/2^k=1$.
We already saw in Example~\ref{exa:longiterone} that $k=\log_2 n$ is the solution.
Therefore, we have
\begin{eqnarray*}
T(n)&=&T(n/2^k)+k\\
&=&T(n/2^{\log_2 n})+\log_2 n\\
&=&T(1)+\log_2 n\\
&=&1+\log_2 n\\
\end{eqnarray*}
Therefore, 
$T(n)=1+\log_2 n$.
\end{answer}
\end{exer}

\newpage

If you can do the following exercise correctly, then you have a firm grasp of the 
iteration method and your algebra skills are superb.  If you have difficulty, keep
working at it and/or get some assistance. I strongly recommend that you do your
best to solve this one on your own.

\begin{exer}
Solve the recurrence relation $T(n)=2T(n-1)+n$, $T(1)=1$.
(Hint:  You will need the result from Exercise~\ref{exer:sumi2i}.)
\vspace*{7.5in}
\begin{answer}
The final answer is $T(n)=2^{n+1}-n-2$ or $T(n)=4\cdot2^{n-1}-n-2$.
It is important that you can work this out yourself, so try your best to get this
answer without looking further.  But if you get stuck or have a different answer,
you can refer to the following skeleton of steps--we have omitted many of the steps
because we want you to work them out.  It does provide a few reference points
along the way, however.
\begin{eqnarray*}
T(n)&=&2T(n-1)+n\\
&=&2(2T(n-2)+(n-1))+n\, \, \, \mbox{(having $n$ instead of $(n-1)$ is a common error)}\\
&=&2^2T(n-2)+3n-2\, \, \, \mbox{(it is unclear yet if I should have $3n-2$ or some other form)}\\
&\vdots& \, \, \, \, \, \, \mbox{(many skipped steps)}\\
&=&2^kT(n-k)+(2^k-1)n-\dsum_{i=1}^{k-1}i2^i\, \, \, \, \mbox{(the all-important pattern revealed)}\\
&\vdots& \, \, \, \, \, \, \mbox{(plug in appropriate value of $k$ and simplify)}\\
&=&2^{n+1}-n-2.
\end{eqnarray*}
\end{answer}
\end{exer}

\newpage

\subsection{Master Theorem} 
We will omit the proof of the following theorem which is particularly
useful for solving recurrence relations that result from the analysis of
certain types of recursive algorithms--especially divide-and-conquer algorithms.

\begin{thm}[Master Theorem]\index{Master Theorem}
\index{recurrence relations!solving!Master Theorem} 
Let $T(n)$ be a monotonically increasing function satisfying 
\begin{eqnarray*}
T(n)&=&aT(n/b)+f(n)\\
T(1)&=&c
\end{eqnarray*}
where $a\geq 1$, $b>1$, and $c>0$.
If $f(n)=\Theta(n^d)$, where $d\geq 0$, then
$$T(n)=\left\{
\begin{array}{ll}
\Theta(n^d)&\text{if } a<b^d\\
\Theta(n^d\log n)&\text{if } a=b^d\\
\Theta(n^{\log_b{a}})&\text{if } a>b^d
\end{array}
\right.$$
\end{thm}

\begin{exa}
Use the Master Theorem to solve the recurrence 
$$T(n) = 4T(n/2) + n, T(1)=1.$$
 \begin{sol} 
We have $a=4, b=2$, and $d = 1$. 
Since $4>2^1$, $T(n)=\Theta(n^{\log_2{4}})=\Theta(n^2)$ by the third case of the Master Theorem.
\end{sol}
\end{exa}

\begin{exa}
Use the Master Theorem to solve the recurrence 
$$T(n) = 4T(n/2) + n^{2}, T(1)=1.$$
\begin{sol} 
We have $a=4, b=2$, and $d=2$.
Since $4=2^2$, we have $T(n)=\Theta(n^2\log n)$ by the second case of the Master Theorem.
\end{sol}
\end{exa}

\begin{exa}
Use the Master Theorem to solve the recurrence 
$$T(n) = 4T(n/2) + n^{3}, T(1)=1.$$
\begin{sol}
Here, $a=4, b=2$, and $d=3$.
Since $4<2^3$, we have $T(n)=\Theta(n^3)$ by the first case of the Master Theorem.
\end{sol}
\end{exa}

Wow.  That was easy.\footnote{Almost {\em too} easy.}
But the ease of use of the Master Method comes with a cost.  Well, two actually.
First, notice that we do not
get an {\em exact} solution, but only an {\em asymptotic bound} on the solution.  Depending on
the context, this may be good enough.  If you need an exact numerical solution, the Master
Method will do you no good.  But when analyzing algorithms, typically we are more interested
in the asymptotic behavior.  In that case, it works great.
Second, it only works for recurrences that have the exact form $T(n)=aT(n/b)+f(n)$.
It won't even work on similar recurrences, 
such as $T(n)=T(n/b)+T(n/c)+f(n)$.


\begin{exer}
Use the Master Theorem to solve the recurrence 
$$T(n) = 2T(n/2) + 1, T(1)=1.$$
\vspace*{1.25in}
\begin{answer}
Here, $a=2, b=2$, and $d=0$. ($d=0$ since $1=1\cdot n^0$.  In general, $c=c\cdot n^0$, so when 
$f(n)$ is a constant, $d=0$.)
Since $a>2^0$, we have $T(n)=\Theta(n^{\log_aa})=\Theta(n^1)=\Theta(n)$ 
by the third case of the Master Theorem.
\end{answer}
\end{exer}

\begin{exa}
Let's redo one from a previous section.  
Use the Master Theorem to solve the recurrence 
$$
R(n) = \left\{ 
\begin{array}{ll}
1&\mbox{when $n = 1$}\\
2 R(n/2) + n/2 & \mbox{otherwise}
\end{array}
\right.
$$
\begin{sol}
Here, we have $a=2$, $b=2$, and $d=1$.  Since $2=2^1$, $R(n)=\Theta(n^1\log n)=\Theta(n\log n)$.
Recall that in Example~\ref{exa:R} we showed that $R(n)=n + (\log_2 n) n/2$. Since
$n + (\log_2 n) n/2=\Theta(n\log n)$, our solution is consistent.
\end{sol}
\end{exa}

\begin{exer}
Use the Master Theorem to solve the recurrence 
$$T(n)=7T(n/2)+15n^2/4, T(1)=1.$$
\vspace*{1.25in}
\begin{answer}
We have $a=7$, $b=2$, and $d=2$.  Since $7>2^2$, the third case of the
Master Theorem applies so $T(n)=\Theta(n^{\log_{2} 7})$, which is about $\Theta(n^{2.8})$.
\end{answer}
\end{exer}

\begin{ques}
In the solution to the previous exercise, we stated that
\begin{quote}
`$T(n)=\Theta(n^{\log_{2} 7})$, which is about $\Theta(n^{2.8})$.'
\end{quote}
Why didn't we just say 
`$T(n)=\Theta(n^{\log_{2} 7})=\Theta(n^{2.8})$'?

\noindent
\answerlinetwo
\begin{answer} 
Because it isn't true.  Although the growth rate of $n^{\log_{2} 7}$ and $n^{2.8}$ are
close, they are not exactly the same, so $\Theta(n^{\log_{2} 7})\neq \Theta(n^{2.8})$.
We {\em could} say that $T(n)=\Theta(n^{\log_{2} 7})=O(n^{2.81})$, but then we have lost
the `tightness' of the bound.  And I want to be able to say ``Yo dawg, that bound is really {\em tight}!''
\end{answer}
\end{ques}

\begin{exer}
We saw in Example~\ref{exa:binSearchRecRel} that the complexity of binary search is
given by the recurrence relation $T(n)=T(n/2)+1$, $T(0)=1$ (and you may assume that $T(1)=1$). 
Use the Master Theorem to solve this recurrence.

\vspace*{1.5in}
\begin{answer} 
Here we have $a=1$, $b=2$, and $d=0$.  Since $1=2^0$,
the second case of the Master Theorem tells is that $T(n)=\Theta(n^0\log n)=\Theta(\log n)$.
Since we have already seen several times that $T(n)=\log_2 n +1$, we can notice that 
this answer is consistent with those.  It's a good thing.
\end{answer}
\end{exer}

\subsection{Linear Recurrence Relations}\label{subsec:linear}

Although in my mind linear recurrence relations are of the least importance of
these four methods for computer scientists, we will discuss them very briefly, both for
completeness sake, and because we can talk about the Fibonacci numbers again.

\index{recurrence relations!solving!linear} 

\begin{df}
Let $c_1, c_2, \ldots , c_k$ be real constants and $f:\BBN \to \BBR$
a function. A recurrence relation of the form 
\begin{equation}\label{eq:linRec}
a_{n} = c_{1}a_{n-1}+c_{2}a_{n-2}+\cdots+c_{k}a_{n-k}+f(n)
\end{equation}
is called a {\bf linear recurrence relation} (or {\bf linear difference equation}). 
\index{linear recurrence relation}
\index{recurrence relations!linear}
If $f(n)=0$ (that is, there is no non-recursive term), 
we say that the equation is {\bf homogeneous}, and otherwise
we say the equation is {\bf nonhomogeneous}.
\index{homogeneous recurrence relation}
\index{nonhomogeneous recurrence relation}
\index{recurrence relation!homogeneous}
\index{recurrence relation!nonhomogeneous}
\end{df}

The {\em order} of the recurrence is  the difference between the
highest and the lowest subscripts. 
\begin{exa} $u_{n} = u_{n -1} + 2$ is of the first order, and
$u_{n} = 9u_{n-4} + n^5$ is of the fourth order.
\end{exa}

There is a general technique that can be used to solve linear homogeneous recurrence relations.
However, we will restrict our discussion to certain first and second order recurrences.

\subsubsection{First Order Recurrences}
\index{first order recurrence}
\index{recurrence relations!first-order} 
\index{recurrence relations!solving!first-order} 

In this section we will learn a technique to solve some first-order recurrences.
We won't go into detail about why the technique works.

\begin{proc}
Let $f(n)$ be a polynomial and $a\neq 1$.  Then the following 
technique can be used to solve a first order linear recurrence
relations of the form
$$x_n = ax_{n - 1} + f(n).$$
\begin{enumerate}
\item First, ignore $f(n)$.  That is, solve the homogeneous recurrence $x_n = ax_{n - 1}$.
This is done as follows:
\begin{lenum}
\item
`Raise the subscripts', so  $x_n = ax_{n - 1}$ becomes $x^n = ax^{n - 1}$. 
This is called the {\em characteristic equation}. \index{characteristic equation}
\item
Canceling this gives $x = a.$ 
\item
The solution to the homogeneous equation $x_n = ax_{n - 1}$
will be of the form $x_n = Aa^n$, where $A$ is a constant to be determined.
\end{lenum} 
\item
Assume that the solution to the original recurrence relation, 
$x_n = ax_{n - 1} + f(n)$, is of the form $x_n = Aa^n + g(n),$ 
where $g$ is a polynomial of the same degree as $f(n)$.
\item
Plug in enough values to determine the correct constants for the coefficients of $g(n)$.
\end{enumerate}
\end{proc}

This procedure is a bit abstract, so let's just jump into seeing it in action.

\begin{exa} 
Let $x_0 = 7$ and $x_n = 2x_{n - 1}, n \geq 1.$ Find a closed form for $x_n.$
\begin{sol} 
Raising subscripts we have the characteristic equation
$x^n = 2x^{n - 1}$. Canceling, $x = 2$. Thus we try a solution of
the form $x_n = A2^n$, were $A$ is a constant. But $7 = x_0 =
A2^0=A$ and so $A = 7.$
The solution is thus $x_n = 7(2)^n$.
% Skip it for now.  What's the point?
% CAC, 6/16/14
%For kicks, here is a different method that sometimes works. 
%We have:
%$$
%\begin{array}{lcl}
%x_0 & = & 7 \\
%x_1 & = & 2x_0 \\
%x_2 & = & 2x_1 \\
%x_3 & = & 2x_2 \\
%\vdots & \vdots & \vdots \\
%x_n & = & 2x_{n-1} \\
%\end{array}
%$$
%Multiplying both columns,
%$$x_0x_1\cdots x_n = 7\cdot 2^nx_0x_1x_2\cdots x_{n - 1}.$$
%Canceling the common factors on both sides of
%the equality,
%$$x_n = 7\cdot 2^n.$$
\end{sol}
\end{exa}

\begin{exa}\label{exa:partSol}
Let $x_0 = 7$ and $x_n = 2x_{n - 1} + 1, n \geq 1.$ Find a closed form for $x_n.$
\begin{sol}
By raising the subscripts in the homogeneous equation we
obtain $x^n = 2x^{n - 1}$ or $x = 2$. A solution to the
homogeneous equation will be of the form $x_n = A(2)^n$. 
Now $f(n)= 1$ is a polynomial of degree 0 (a constant) and 
so the general solution should have
the form $x_n = A2^n + B$. Now, $7 = x_0 = A2^0 + B = A + B$.
Also, $x_1 = 2x_0 + 1 = 15$ and so $15 = x_1 = 2A + B$. Solving
the simultaneous equations
$$A + B = 7,$$
$$2A + B = 15,$$
Using these equations, we can see that $A=7-B$ and $B=15-2A$.
Plugging the latter into the former, we have $A=7-(15-2A)=-8+2A$, or $A=8$.
Plugging this back into either equation, we can see that $B=-1$.
So the solution is $x_n = 8(2^n) - 1 = 2^{n + 3} - 1.$
\end{sol}
\end{exa}

\newpage
\begin{exer}
Let $x_0 = 2, x_n = 9x_{n - 1} - 56n + 63$. Find a closed form for
this recurrence.
\vspace*{3.75in}
\begin{answer}
By raising the subscripts in the homogeneous equation we
obtain the characteristic equation $x^n = 9x^{n - 1}$ or $x = 9$.
A solution to the homogeneous equation will be of the form $x_n =
A(9)^n$. Now $f(n) = -56n + 63$ is a polynomial of degree 1 and so
we assume that the solution will have the form $x_n = A9^n + Bn + C$. 
Now $x_0 = 2,
x_1 = 9(2) - 56 + 63 = 25, x_2 = 9(25) - 56(2) + 63 = 176$. We
thus solve the system
$$2 = A + C,$$
$$25 = 9A + B + C,$$
$$176 = 81A + 2B + C.$$
We find $A = 2, B = 7, C = 0$, so the solution is $x_n = 2(9^n) + 7n.$
\end{answer}
\end{exer}


% Too many variables to solve for.
% CAC, 6/16/14
%\begin{exa}
%Let $x_0 = 1, x_n = 3x_{n - 1} - 2n^2 + 6n - 3$. Find a closed
%form for this recursion.
%\begin{sol}  By raising the subscripts in the homogeneous equation we
%obtain the characteristic equation $x^n = 3x^{n - 1}$ or $x = 9$.
%A solution to the homogeneous equation will be of the form $x_n =
%A(3)^n$. Now $f(n) = - 2n^2 + 6n - 3$ is a polynomial of degree 2
%and so we test a particular solution of the form $Bn^2 + Cn + D$.
%The general solution will have the form $x_n = A3^n + Bn^2 + Cn +
%D$. Now $x_0 = 1, x_1 = 3(1) - 2 + 6 - 3 = 4, x_2 = 3(4) - 2(2)^2
%+ 6(2) - 3 = 13, x_3 = 3(13) - 2(3)^2 + 6(3) - 3 = 36$. 
%We thus solve the system
%$$1 = A + D,$$
%$$4 = 3A + B + C + D,$$
%$$13 = 9A + 4B + 2C + D,$$
%$$36 = 27A + 9B + 3C + D.$$
%We find $A = B = 1, C = D = 0.$ The general solution is $x_n = 3^n
%+ n^2.$
%\end{sol}
%\end{exa}


% The discussion above mentions $f(n)$ being a polynomial, so why
%an example with an exponential?
%\begin{exa}
%Find a closed form for $x_{n} = 2x_{n - 1} + 3^{n - 1}, x_0 = 2.$
%\begin{sol}  We test a solution of the form $x_n = A2^n + B3^n$. Then
%$x_0 = 2, x_1 = 2(2) + 3^0 = 5.$ We solve the system
%$$2 = A + B,$$
%$$7 = 2A + 3B.$$We find $A = 1, B = 1.$ The general solution is $x_n = 2^n + 3^n.$
%\end{sol}
%\end{exa}

% I don't want to get into this special case.
% CAC, 6/16/14
%
%We now tackle the case when $a = 1.$ In this case, we simply
%consider a polynomial $g$ of degree 1 higher than the degree of
%$f$.
%
%\begin{exa} Let $x_0 = 7$ and $x_n = x_{n - 1} + n, n \geq 1.$ Find a closed formula for $x_n.$ 
%\begin{sol}  
%By raising the subscripts in the homogeneous equation we
%obtain the characteristic equation $x^n = x^{n - 1}$ or $x = 1$. A
%solution to the homogeneous equation will be of the form $x_n =
%A(1)^n = A$, a constant. Now $f(n) = n$ is a polynomial of degree
%1 and so we test a particular solution of the form $Bn^2 + Cn +
%D$, one more degree than that of $f$. The general solution will
%have the form $x_n = A + Bn^2  + Cn + D$. Since $A$ and $D$ are
%constants, we may combine them to obtain $x_n = Bn^2 + Cn + E.$
%Now, $x_0 = 7, x_1 = 7 + 1 = 8, x_2 = 8 + 2 = 10.$ So we solve the
%system
%$$7 = E,$$
%$$8 = B + C + E,$$
%$$10 = 4B + 2C + E.$$ We find $\dis{B = C = \frac{1}{2}, E = 7}$. The general solution is
%$\dis{x_n = \frac{n^2}{2} + \frac{n}{2} + 7}$.
%
%Here is another way to solve it.
% We have
%$$
%\begin{array}{lcl}
%x_0 & = & 7 \\
%x_1 & = & x_0 + 1 \\
%x_2 & = & x_1  + 2\\
%x_3 & = & x_2  + 3\\
%\vdots & \vdots & \vdots \\
%x_n & = & x_{n-1} + n \\
%\end{array}
%$$
%Adding both columns,
%$$x_0 + x_1 + x_2 + \cdots + x_n = 7 + x_0 + x_2 + \cdots + x_{n - 1} + (1 + 2 + 3 + \cdots + n).$$
%Cancelling and using the fact that $\dis{1 + 2 + \cdots + n =
%\frac{n(n + 1)}{2}}$,
%$$x_n = 7 + \frac{n(n + 1)}{2}.$$
%\end{sol}
%\end{exa}

\subsubsection{Second Order Recurrences} 
Let us now briefly examine how to solve some
second order recursions.
\index{second order recurrence}
\index{recurrence relations!second-order} 
\index{recurrence relations!solving!second-order} 

\begin{proc} 
Here is how to solve a second-order homogeneous
linear recurrence relations of the form
$$x_n = ax_{n - 1} + bx_{n - 2}.$$
\begin{enumerate}
\item Find the characteristic equation  by ``raising the
subscripts.'' We obtain $x^n = ax^{n - 1}  + bx^{n - 2}$.
\item
Canceling this gives $x^2 - ax - b = 0.$ This equation has two
roots $r_1$ and $r_2.$
\index{characteristic equation}
\item If the roots are different, the
solution will be of the form $x_n = A(r_1)^n + B(r_2)^n$, where
$A, B$ are  constants. \item If the roots are identical, the
solution will be of the form $x_n = A(r_1)^n + Bn(r_1)^n$.
\end{enumerate}
\end{proc}

\begin{exa}
Let $x_0 = 1, x_1 = -1, x_{n + 2} + 5x_{n + 1} + 6x_n = 0$.
\begin{sol}  The characteristic equation is $x^2 + 5x + 6 = (x + 3)(x
+ 2) = 0$. Thus we test a solution of the form $x_n = A(-2)^n +
B(-3)^n.$ Since $1 = x_0 = A + B$, and $-1 = -2A - 3B$, we quickly find
$A = 2$, and $B = -1.$ Thus the solution is $x_{n} = 2(-2)^n -(-3)^n.$
\end{sol}
\end{exa}

\begin{exa}\label{exa:fibSol}\index{Fibonacci numbers}
Find a closed form for the Fibonacci recurrence $f_0 = 0, f_1 = 1,
f_n = f_{n - 1} + f_{n - 2}$.
\begin{sol}  The characteristic equation is $f^2 - f - 1 = 0$.
This has roots $\frac{1 \pm \sqrt{5}}{2}$.
Therefore, a solution will have the form $${f_n = A\left(\frac{1 +
\sqrt{5}}{2}\right)^n + B\left(\frac{1 - \sqrt{5}}{2}\right)^n}.$$
The initial conditions give
$$0 = A + B, \mbox{ and}$$
$$1 = A\left(\frac{1 + \sqrt{5}}{2}\right) + B\left(\frac{1 - \sqrt{5}}{2}\right)
= \frac{1}{2}\left(A + B\right) + \frac{\sqrt{5}}{2}\left(A -
B\right) = \frac{\sqrt{5}}{2}\left(A - B\right).$$ 
From these two equations, we obtain
$\dis{A = \frac{1}{\sqrt{5}}, B = -\frac{1}{\sqrt{5}}}$. 
We thus have 
$$
f_n = \frac{1}{\sqrt{5}}\left(\frac{1 + \sqrt{5}}{2}\right)^n -
\frac{1}{\sqrt{5}}\left(\frac{1 - \sqrt{5}}{2}\right)^n.
$$
\end{sol}
\end{exa}

%\begin{exa}
%Why was the recursive algorithm to compute $f_n$ from Example \ref{exa:fibAlg} so bad?
%\end{exa}
%\begin{sol}
%Let's count the number of additions \scode{FibR(n)} computes since that is the main
%thing that the algorithm does.
%Let $F(n)$ be the number of additions required to compute $f_n$ using \scode{FibR(n)}.
%Since \scode{FibR(n)} performs one addition and then calls \scode{FibR(n-1)} and \scode{FibR(n-2)},
%it is easy to see that 
%$$F(n)=F(n-1)+F(n-2)+1,$$
%where $F(0)=F(1)=0$ is clear from the algorithm.
%%From Example \ref{exa:fibSol} and following the ideas from Example \label{exa:partSol},
%%we know that $$F(n)= A\frac{1}{\sqrt{5}}\left(\frac{1 + \sqrt{5}}{2}\right)^n -
%%B\frac{1}{\sqrt{5}}\left(\frac{1 - \sqrt{5}}{2}\right)^n+C$$
%%for some constants $A$, $B$, and $C$.
%%From the algorithm, it is clear that $F(0)=F(1)=0$, and it is easy to compute $F(2)=F(1)+F(0)+1 = 1$.
%%Therefore 
%%\begin{eqnarray*}
%%0=F(0)&=&A\frac{1}{\sqrt{5}}\left(\frac{1 + \sqrt{5}}{2}\right)^0 +
%%B\frac{1}{\sqrt{5}}\left(\frac{1 - \sqrt{5}}{2}\right)^0+C
%%=A\frac{1}{\sqrt{5}}+B\frac{1}{\sqrt{5}}+C\\
%%0=F(1)&=&A\frac{1}{\sqrt{5}}\left(\frac{1 + \sqrt{5}}{2}\right) +
%%B\frac{1}{\sqrt{5}}\left(\frac{1 - \sqrt{5}}{2}\right)+C\\
%%1=F(2)&=&A\frac{1}{\sqrt{5}}\left(\frac{1 + \sqrt{5}}{2}\right)^2 +
%%B\frac{1}{\sqrt{5}}\left(\frac{1 - \sqrt{5}}{2}\right)^2+C\\
%%\end{eqnarray*}
%\end{sol}

\begin{exer}
Find a closed form for the recurrence $x_0 = 1, x_1 = 4, x_n = 4x_{n - 1}  - 4x_{n -
2}.$
\vspace*{4in}
\begin{answer} 
The characteristic equation is $x^2 - 4x + 4 = (x - 2)^2
= 0$. There is a multiple root and so we must test a solution of
the form $x_n = A2^n + Bn2^n.$ The initial conditions give
$$1 = A,$$
$$4 = 2A + 2B.$$ 
This solves to $A = 1, B = 1.$ The solution is thus $x_n = 2^n + n2^n$.
\end{answer}
\end{exer}


% Interesting example, but what point does it serve?
%CAC, 6/12/14
%\begin{exer}
%Solve the recursion $a_n=1+\sum _{k=1} ^{n-1} a_k$ for $n\geq 2$ and  $a_1=1$.
%\begin{answer}
%Observe that
%$$a_n-a_{n-1} =\left(1+\sum _{k=1} ^{n-1} a_k\right)-\left(1+\sum _{k=1} ^{n-2} a_k\right)=a_{n-1}.   $$
%This means that $a_n=2a_{n-1}$ and so
%$$ \begin{array}{lll} a_n& = & 2a_{n-1}\\  a_{n-1}& = & 2a_{n-2}\\     \vdots & \vdots & \vdots \\  a_2& = & 2a_1\\ \end{array}$$
%Multiplying all these equalities,
%$$a_na_{n-1}\cdots a_2 = 2^{n-1}a_{n-1}a_{n-2}\cdots a_1 \implies a_n = 2^{n-1}a_1 = 2^{n-1}. $$
%\end{answer}
%\end{exer}

% Also interesting, but perhaps beyond the scope/purpose of this chapter.
%CAC, 6/12/14
%\begin{exer}[Derangements] An absent-minded secretary is filling $n$
%envelopes with $n$ letters. Find a recurrence relation for the number $D_n$ of
%ways in which she never stuffs the right letter into the right
%envelope. 
%\begin{answer} Number the envelopes $1, 2, 3, \cdots , n$.
%We condition on the last envelope.  Two events might happen. Either
%$n$ and $r$ (for some $1 \leq r \leq n - 1)$ trade places or they do not.
%
% In the first case, the two letters $r$ and $n$ are
%misplaced. Our task is just to misplace the other $n - 2$ letters,
%$(1, 2, \cdots , r - 1, r + 1, \cdots , n - 1)$ in the slots $(1, 2,
%\cdots , r - 1, r + 1, \cdots , n - 1).$ This can be done in $D_{n -
%2}$ ways. Since $r$ can
%be chosen in $n - 1$ ways, the first case can happen in $(n - 1)D_{n - 2}$ ways.
%
%In the second case, let us say that letter $r$, $(1 \leq r \leq n -
%1)$ moves to the $n$-th position but $n$ moves not to the $r$-th
%position.  Since $r$ has been misplaced, we can just ignore it.
%Since $n$ is not going to the $r$-th position, we may relabel $n$ as
%$r$.  We now have $n - 1$ numbers to misplace, and this can be done
%in
%$D_{n - 1}$ ways.
%
%As $r$ can be chosen in $n - 1$ ways, the total number of ways for
%the second case is $(n - 1)D_{n - 1}.$  Thus $D_n = (n - 1)D_{n - 2}
%+ (n - 1)D_{n - 1}.$
%\end{answer}
%\end{exer}

\newpage

\section{Analyzing Recursive Algorithms}\label{sec:anaRecAlg}
In Section \ref{sec:recRel} we already saw a few examples of analyzing recursive algorithms.
We will provide a few more examples in this section. 
In case it isn't clear, the most common method to analyze a recursive algorithm is to develop
and solve a recurrence relation for its running time.  Let's see some examples.

The two most two popular recursive sorting algorithms
are {\em merge sort} and {\em Quicksort}. We will briefly discuss
each algorithm and then analyze the running time of each.

The basic idea behind {\em merge sort} is simple:\index{merge sort}\index{sort!merge}
Divide an array in half. Sort the first half, then sort the
second half, then merge the two halves together.
The first question you might be asking is how do you sort the
two halves. Easy: merge sort them! But we need a base case
otherwise the recursion will go on forever. That is also easy:
An array of size 0 or 1 is clearly already sorted.

Now we are left with two questions: how do we merge the two 
halves and how do we use that to implement the entire algorithm?
Here we will simply provide both algorithms without much
explanation and we will focus on analyzing them. If you want to
better understand both merge sort and Quicksort, I recommend you 
take a data structures and algorithms course.

\begin{exa}\index{merge}
The \verb+Merge+ algorithm needs to take two sorted arrays
and combine them into a single sorted array.
More accurately, it takes as input an array $A$ which
is sorted from $left$ to $mid$ and from $mid+1$ to $right$,
and modifies $A$ so it is sorted from $left$ to $right$.

For instance, we might have 
$left=5$, $right=15$, $mid=9$, and 
$$A=[73, 23, 45, 67, 34, \underline{23, 37, 41, 45, 65}, \overline{12, 28, 39, 59, 68, 89}, 12, 34, 48, 47],$$
where the underlined portion, indexed from $left$ to $mid$, is sorted,
and the overlined portion, indexed from $mid+1$ to $right$ is also sorted,
and we want merge to return the array
$$A=[73, 23, 45, 67, 34, \underline{12, 23, 28, 37, 39, 41, 45, 59,  65, 68, 89}, 12, 34, 48, 47]$$
in which the underlined portion, indexed from $left$ to $right$
is now fully sorted (check carefully that it is correct).

Now that we better understand what \verb+Merge+ is supposed to do,
we can present an algorithm that accomplishes the task.
\begin{code}
void Merge(int A[],int left,int mid,int right) {
    int n=right-left+1;
    int m=mid-left+1;
    int *B=new int[n]; // Create new array
    for(int i=0;i<m;i++)
        B[i]=A[left+i]; // copy 1st half forward
    for(int j=m;j<n;j++)
        B[j]=A[right-j+m]; // copy 2nd half reversed
    int i=0; 
    int j=n-1;
    for(int k=left;k<=right;k++) // Now merge back to A
        if(B[i]<B[j]) 
            A[k]=B[i++];
        else
            A[k]=B[j--];
}
\end{code}
If we set $n=right-left+1$, the size of the portion of the array that we are dealing with,
we should be able to easily see that \verb+Merge+ takes $O(n)$ time since 
it executes two loops whose total number of iterations is $n$, and the third loop
which has $n$ iterations, and the rest of the code takes constant time.
\end{exa}

\begin{exa}
Now that we have the \verb+Merge+ algorithm, implementing merge sort is pretty straightforward 
given the description above. 

\begin{code}\index{merge sort}\index{sort!merge}
MergeSort(int[] A,int left,int right) {
   if (left < right) {
       int mid = (left + right)/2;
       MergeSort(A, left, mid);
       MergeSort(A, mid + 1, right);
       Merge(A, left, mid, right);
   }
}
\end{code}

Notice that if $left\geq right$ the algorithm does nothing.
Also notice that this corresponds to subarrays with 0 or 1 elements, which 
clear do not need to be sorted. This is the base case we mentioned above.
\end{exa}

\begin{exa}\label{exa:meregesortAnal}
\index{merge sort}
What is the worst-case running time of merge sort?
\begin{sol}
The algorithm for \verb+MergeSort+ is below.
Let $T(n)$ be the worst-case running time of 
\verb+MergeSort+ on an array of size $n=right-left$.
Recall that \verb+Merge+ takes two sorted arrays and merges them into
one sorted array in time $\Theta(n)$, where $n$ is the number of
elements in both arrays.\footnote{Since our goal here is to analyze the algorithm,
we won't provide a detailed implementation of \texttt{Merge}. % \verb doesn't work in footnotes?!  
All we need to know is its complexity.}  
Since the two recursive calls to \verb+MergeSort+
 are on arrays
of half the size, they each require time $T(n/2)$ in the worst-case.
The other operations take constant time.  Below we annotate the \verb+MergeSort+ algorithm
with these running times.
\begin{center}
\begin{tabular}{|l|c|}
\hline
Algorithm&Time required\\
\hline
\verb-MergeSort(int[] A,int left,int right) {-&$T(n)$\\
\verb-   if (left < right) {-                 &$C_1$\\
\verb-       int mid = (left + right)/2;-     &$C_2$\\
\verb-       MergeSort(A, left, mid);-        &$T(n/2)$\\
\verb-       MergeSort(A, mid + 1, right);-   &$T(n/2)$\\
\verb-       Merge(A, left, mid, right);-     &$\Theta(n)\leq C_3n$\\
\verb-   }-&\\
\verb-}-&\\
\hline
\end{tabular}
\end{center}

Given this, we can see that
\begin{eqnarray*}
T(n)&=&C_1+C_2+T(n/2)+T(n/2)+\Theta(n)\\
&=&2 T(n/2) +\Theta(n).
\end{eqnarray*}
Notice that we absorbed the constants $C_1$ and $C_2$ into the $\Theta(n)$ term.
For simplicity, we will also replace the $\Theta(n)$ term with $cn$ (where $c$ is a constant)
and rewrite this as
 $$T(n)=2 T(n/2)+c n.$$
We could use the Master Theorem to prove that $T(n)=\Theta(n\log n)$,
 but that would be too easy. 
Instead, we will use induction to prove that $T(n)=O(n\log n)$, 
and leave the $\Omega$-bound to the reader.

By definition, $T(n)=O(n\log n)$ if and only if there exists constants
$k$ and $n_0$ such that $T(n)\leq k n\log n$ for all $n\geq n_0$.

For the base case, notice that $T(2)=a$ for some constant $a$,
and $a\leq k 2\log 2= 2 k$ as long as we pick $k\geq a/2$.
Now, assume that $T(n/2)\leq k(n/2)\log (n/2)$. Then

\begin{eqnarray*}
T(n)&=&2 T(n/2) + c n\\
&\leq&2(k(n/2)\log(n/2)+ c n\\
&=&k n \log(n/2)+ c n\\
&=&k n \log n - k n\log 2 + c n\\
&=&k n\log n + (c-k)n\\
&\leq& k n\log n \,\,\,\,\, \mbox{ if } k\geq c
\end{eqnarray*}
As long as we pick $k=\max\{a/2,c\}$, we have 
$T(n)\leq k n\log n$, so $T(n)=O(n\log n)$ as desired.
\end{sol}
\end{exa}

\begin{exer}
We stated in the previous example that we could use the
Master Theorem to prove that if $T(n)=2 T(n/2)+c n,$
then $T(n)=\Theta(n\log n)$.  Verify this.
\vspace*{1.5in}
\begin{answer}
We have $a=2$, $b=2$, and $d=1$. Since $2=2^1$, we have that 
$T(n)=\Theta(n\log n)$ by the second case of the Master Theorem.
\end{answer}
\end{exer}

\begin{ques}
Answer  the following questions about points that were made in Example~\ref{exa:meregesortAnal}.
\begin{lenum}
\item
Why were we allowed to absorb the constants $C_1$ and $C_2$ into the $\Theta(n)$ term?

\noindent
\answerlinethree
\item
Why were we able to replace the $\Theta(n)$ term with $cn$?

\noindent
\answerlinethree
\end{lenum}
\begin{answer}
(a) $C_1$ and $C_2$ are of lower order than $n$.
Thus, we can do this according to part (c) of Theorem~\ref{thm:asyprops}.
(b) We know that the $\Theta(n)$ represents some function $f(n)$.  By definition of $\Theta$,
there are constants, $c_1$ and $c_2$ such that $c_1n\leq f(n)\leq c_2n$ for all $n\geq n_0$
for some constant $n_0$.  So we essentially replaced $f(n)$ with $c_2n$ (since we are
looking for a worst-case). 
Note that it doesn't matter what $c_2$ is.  
We know there is some constant that works, so we just call it $c$. 
\end{answer}
\end{ques}


\begin{exa}[Towers of Hanoi]\index{Towers of Hanoi}\index{Hanoi, Towers of}
The following legend is attributed to
 French mathematician Edouard Lucas in 1883. 
In an Indian temple there are 64 gold disks resting on three pegs. 
At the beginning of time, God placed these disks on the first peg
and ordained that a group of priests should transfer them to the third peg according to the following rules:
\begin{enumerate}
\item The disks are initially stacked on peg A, in decreasing order (from bottom to top).
\item The disks must be moved to another peg in such a way that only one disk 
is moved at a time and without stacking a larger disk onto a smaller disk.
\end{enumerate}
When they finish, the Tower will crumble and the world will end. 
How many moves does it take to solve the {\em Towers of Hanoi} problem with $n$ disks?
\begin{sol}
The usual (and best) algorithm to solve the 
{\em Towers of Hanoi} is:
\begin{itemize}
\item
Move the top $n-1$ disk to from peg 1 to peg 2.
\item
Move the last disk from peg 1 to peg 3.
\item
Move the top $n-1$ disks from peg 2 to peg 3.
\end{itemize}
The only question is how to move the top $n-1$ disks.
The answer is simple: use recursion but switch the peg numbers.
Here is an implementation of this idea:
\begin{code}
  void solveHanoi(int N,int source,int dest,int spare) {
      if(N==1) {
          moveDisk(source,dest);
      } else {
          solveHanoi(N-1,source,spare,dest);
          moveDisk(source,dest);
          solveHanoi(N-1,spare,dest,source);
      }
  }
\end{code}
Don't worry if you don't see why this algorithm works.
Our main concern here is analyzing the algorithm.

The exact details of \verb+moveDisk+ depend on how the pegs/disks
are implemented, so we won't provide an implementation of it.
But it doesn't actually matter anyway since we just need to count 
the number of times \verb+moveDisk+ is called.  As it turns out,
any reasonable implementation of \verb+moveDisk+ will take
constant time, so the complexity of the algorithm is essentially
the same as the number of calls to \verb+moveDisk+.

Let $H(n)$ be the number of moves it takes to solve the {\em Towers of Hanoi} problem 
with $n$ disks.  
Then $H(n)$ is the number of times \verb+moveDisk+ is called
when running \verb+solveHanoi(n,1,2,3)+.
It should be clear that $H(1)=1$ since the algorithm simply makes a single
call to \verb+moveDisk+ and quits.
When $n>1$, the algorithm makes two calls to \verb+solveHanoi+ with the
first parameter being $n-1$ and one call to \verb+moveDisk+.
Therefore, we can see that

\noindent
\hspace*{2in} $H(n)=2H(n-1)+1.$

\noindent
As with the first example, we want a closed form for $H(n)$.
But we already showed that $H(n)=2^{n}-1$ in Examples \ref{exa:Hanoi1} and \ref{exa:Hanoi2}.
\end{sol}
\end{exa}

\begin{exer}
Let $T(n)$  be the complexity of \verb+blarg(n)+. 
Give a recurrence relation for $T(n)$.
\begin{code}
int blarg(int n) {
     if(n>5) {
        return blarg(n-1)+blarg(n-1)+blarg(n-5)+blarg(sqrt(n));
     }
     else {
        return n;
     }
}
\end{code}
\answerline
\begin{answer}
$T(n)=2 T(n-1)+T(n-5)+T(\sqrt{n})+1$ if $n>5$,
$T(1)=T(2)=T(3)=T(4)=T(5)=1$.
You can also have $+c$ instead of $+1$ in the recursive definition.
\end{answer}
\end{exer}

\begin{exer}\label{exer:stooge}\index{Stooge sort}\index{sort!Stooge}
Give a recurrence relation for the running time of \verb+stoogeSort(A,0,n-1)+.
(Hint: Start by letting $T(n)$ be the running time of \verb+stoogeSort+ on an array of size $n$.)
\begin{code}
  void stoogeSort(int[] A,int L,int R){
      if(R<=L) return; // Array has at most one element
      if(A[R]<A[L]) {  // Swap first and last element 
          Swap(A,L,R); // if they are out of order
      }
      if(R-L>1){ // If the list has at least 2 elements
          int third=(R-L+1)/3;
          stoogeSort(A,L,R-third); // Sort first two-thirds
          stoogeSort(A,L+third,R); // Sort last two-thirds
          stoogeSort(A,L,R-third); // Sort first two-thirds again
      }
  }
\end{code}
\answerline
\begin{answer}
Beyond the recursive calls, \verb+StoogeSort+ does only a constant amount of work--we'll call it $1$.
Then it makes three calls with sub-arrays of size $(2/3)n$.  Therefore, 
$T(n)=3T((2/3)n)+1$ or $T(n)=3T(2n/3)+1$, with base case $T(1)=1$.
\end{answer}
\end{exer}

\newpage

\begin{exer}
Solve the recurrence relation you developed for \verb+StoogeSort+ in the previous exercise.
(Make sure you verify your solution to the previous problem before you attempt to solve
your recurrence relation).
\vspace*{2in}
\begin{answer}
We can use the Master Theorem for this one.  
$a=3$, $b=3/2$ and $d=0$.  
(Notice that $b\neq 2/3$!  If you made this mistake, make sure you understand why it is incorrect.)
Since $3>1=(3/2)^0$, the third case of the Master Theorem tells us that
$T(n)=\Theta\left(n^{\log_{3/2}(3)}\right)$.  Although this looks weird, we can have a rational number as
the base of a logarithm (in fact, the base of $ln(n)$ is $e$, an irrational number).
It might be helpful to compute the $\log$ since it isn't clear how good or bad this complexity is.
Notice that $\log_{3/2}(3)\approx 2.71$, so $T(n)$ is approximately $\Theta(n^{2.71})$, but 
$T(n)\neq \Theta(n^{2.71})$, so resist the urge to place an equals sign between these.
\end{answer}
\end{exer}

\begin{ques}
Which sorting algorithm is faster, \verb+MergeSort+ or \verb+StoogeSort+?
Justify your answer.

\noindent
\answerlinethree
\begin{answer}
The complexity of \verb+MergeSort+ is $\Theta(n\log n)$ and the complexity of
\verb+StoogeSort+ is $\Theta\left(n^{\log_\frac{3}{2}(3)}\right)$
which is $\Omega(n^{2.7})$.
Clearly $\Theta\left(n^{\log_\frac{3}{2}(3)}\right)$ grows faster than $\Theta(n\log n)$,
so \verb+MergeSort+ is faster. 
{\em Remember: Faster growth rate means slower algorithm!}
\end{answer}
\end{ques}

\begin{exer}
Give and solve a recurrence relation for the running time of an algorithm that does as follows:
The algorithm is given an input array of size $n$.
If $n<3$, the algorithm does nothing.
If $n\geq 3$, create 5 separate arrays, 
each one-third of the size of the original array.
This takes $\Theta(n)$ to accomplish.  
Then call the same algorithm on each of the 5 arrays.
\vspace*{3in}
\begin{answer}
Let $T(n)$ be the complexity of this algorithm.
From the description, it seems pretty clear that $T(n)=5T(n/3)+cn$.
Using the Master Theorem with $a=5$, $b=3$, and $d=1$, we see that $5>3^1$, 
so the third case applies and $T(n)=\Theta\left(n^{\log_3(5)}\right)$,
which is approximately $\Theta(n^{1.46})$.
\end{answer}
\end{exer}

\subsection{Analyzing Quicksort}\label{sec:qsa}

In this section we give a proof that the average case running time of
randomized quicksort is $\Theta(n\log n)$.\index{Quicksort}\index{sort!Quicksort}
This proof gets its own section because the analysis is fairly involved.
This proof is based on the one presented in Section 8.4 of
the classic {\em Introduction to Algorithms} by Cormen, Leiserson, and Rivest.
The algorithm they give is slightly different, 
and they include some interesting insights, so read their proof/discussion if 
you get a chance.
  
There are several slight variations of
the quicksort algorithm, and 
although the exact running times are different
for each, the asymptotic running times
are all the same. 
Below is the version of \verb+Quicksort+ we will analyze.

\begin{exa}\label{exa:quicksortImp}
Here is one implementation of \verb+Quicksort+:

\noindent
\begin{minipage}[t]{.45\textwidth}
\begin{code}
Quicksort(int A[],int l,int r){
   if (r > l) {
     int p = Partition(A,l,r);
     Quicksort(A,l,p-1);
     Quicksort(A,p+1,r);
   }
}
\end{code}
\end{minipage}
\hfill
\begin{minipage}[t]{.52\textwidth}
\begin{code}
int Partition(int A[],int l,int r){ 
    int piv=l+(rand()%(r-l+1));
    swap(A,l,piv);
    int i = l+1;
    int j = r;
    while (1) {
        while (A[i] <= A[l] && i<r) 
            i++;
        while (A[j] >= A[l] && j>l) 
            j--;
        if (i >= j) {
           swap(A,j,l);
           return j;
         }
        else swap(A,i,j);
    }
}
\end{code}
\end{minipage}
\end{exa}

We will base our analysis on this version of \verb+Quicksort+.
It is straightforward to see that the runtime of 
\verb+Partition+ is $\Theta(n)$ (Problem~\ref{pro:rpartComp}
asks you to prove this). 
We start by developing a recurrence
relation for the average case runtime of \verb+Quicksort+.
\begin{thm}
Let $T(n)$ be the average case runtime of \verb+Quicksort+ on an
array of size $n$.  Then 
$$T(n)=\frac{2}{ n}\sum_{k=1}^{n-1} T(k) +\Theta(n).$$

\noindent
{\bf Proof: }
Since the pivot element is chosen randomly, it is equally likely that
the pivot will end up at any position from $l$ to $r$.
That is, the probability that the pivot ends up at location $l+i$ 
is $1/n$ for each $i=0,\ldots, r-l$.  If we average over all of the possible
pivot locations, we obtain (the last step holds since $T(0)=0$)
\begin{eqnarray*}
T(n)&=&\frac{1}{ n}\left( \sum_{k=0}^{n-1}\left(T(k)+T(n-k-1)\right)\right)
+\Theta(n)\\
&=&\frac{1}{ n}\sum_{k=0}^{n-1} T(k) 
+\frac{1}{ n}\sum_{k=0}^{n-1}T(n-k-1) +\Theta(n)\\
&=&\frac{1}{ n}\sum_{k=0}^{n-1} T(k) 
+\frac{1}{ n}\sum_{k=0}^{n-1}T(k) +\Theta(n)\\
&=&\frac{2}{ n}\sum_{k=0}^{n-1} T(k) +\Theta(n)\\
&=&\frac{2}{ n}\sum_{k=1}^{n-1} T(k) +\Theta(n).
\end{eqnarray*}
\end{thm}

We will need the following result in order to solve the recurrence relation.

\begin{lem}
For any $n\geq 3$, 
$$\sum_{k=2}^{n-1} k\log k \leq \frac{1}{2}n^2\log n - \frac{1}{8}n^2.$$
\begin{pf}
We can write the sum as
$$\sum_{k=2}^{n-1} k\log k =
\sum_{k=2}^{\lceil n/2 \rceil-1} k\log k +
\sum_{k=\lceil n/2 \rceil}^{n-1} k\log k $$
Then we can bound $(k\log k)$ by $(k\log(n/2))=k(\log n - 1)$ in the first sum, 
and by $(k\log n)$ in the second sum.
This gives 

\begin{eqnarray*}
\sum_{k=2}^{n-1} k\log k &=&
\sum_{k=2}^{\lceil n/2 \rceil-1} k\log k + 
\sum_{k=\lceil n/2 \rceil}^{n-1} k\log k\\
&\leq& \sum_{k=2}^{\lceil n/2 \rceil-1} k(\log n -1) + 
\sum_{k=\lceil n/2 \rceil}^{n-1} k\log n\\
&=&(\log n - 1)\sum_{k=2}^{\lceil n/2 \rceil-1} k +
\log n\sum_{k=\lceil n/2 \rceil}^{n-1} k\\
&=&\log n\sum_{k=2}^{\lceil n/2 \rceil-1} k -
\sum_{k=2}^{\lceil n/2 \rceil-1} k +
\log n\sum_{k=\lceil n/2 \rceil}^{n-1} k\\
&=&\log n \sum_{k=2}^{n-1} k - \sum_{k=2}^{\lceil n/2 \rceil-1} k\\
&\leq&\log n \sum_{k=1}^{n-1} k - \sum_{k=1}^{\lceil n/2 \rceil-1} k\\
&\leq&(\log n) \frac{1}{ 2}(n-1)n  - \frac{1}{ 2}(\frac{n}{ 2}-1)\frac{n}{ 2}\\
&=&\frac{1}{ 2}n^2\log n  - \frac{n}{ 2}\log n- \frac{1}{ 8}n^2+ \frac{n}{ 4}\\
&\leq&\frac{1}{ 2}n^2\log n  - \frac{1}{ 8}n^2.\\
\end{eqnarray*}
The last step holds since 
$$\frac{n}{ 4}\leq \frac{n}{ 2}\log n,$$
when $n\geq 3$.
\end{pf}
\end{lem}

\newpage

Now we are ready for the final analysis.

\begin{thm}
Let $T(n)$ be the average case runtime of \verb+Quicksort+ on an
array of size $n$.  Then 
$$T(n)=\Theta(n\log n).$$
\begin{pf}
We need to show that $T(n)=O(n\log n)$ and $T(n)=\Omega(n\log n)$.
To prove that $T(n)=O(n\log n)$, we will show that for some constant
$a$, 
$$T(n)\leq a n\log n \mbox{\ \  for all } n\geq 2.\footnote{We pick 2 for the 
base case since $n\log n$=0 if $n=1$, so we cannot make
the inequality hold. Another solution would be to show
that $T(n)\leq a n\log n + b$. In this case, $b$ can be chosen
so that the inequality holds for $n=1$.}$$
When $n=2$, 
$$a n\log n=a 2 \log 2=2 a,$$
and $a$ can be chosen
large enough so that $T(2)\leq 2a$.  Thus, the inequality holds for the
base case.  
Let $T(1)=C$, for some constant $C$.
For $2<k<n$, assume $T(k)\leq a k \log k$.   Then
$$
\begin{array}{lcll}
T(n)&=&{\displaystyle \frac{2}{n}\sum_{k=1}^{n-1} T(k) +\Theta(n)}\\[.25in]
&\leq&{\displaystyle \frac{2}{n}\sum_{k=2}^{n-1} a k \log k 
+\frac{2}{n}T(1)+\Theta(n)}&
(\mbox{by assumption})\\[.25in]
&=&{\displaystyle \frac{2 a }{ n}\sum_{k=2}^{n-1} k \log k 
+\frac{2}{n} C +\Theta(n)}\\[.25in]
&\leq&{\displaystyle \frac{2 a }{ n}\sum_{k=2}^{n-1} k \log k + C +\Theta(n)}& 
(\mbox{since } \frac{2 }{ n}\leq 1)\\[.25in]
&\leq&{\displaystyle \frac{2 a }{ n}\left( \frac{1}{2}n^2\log n - \frac{1}{8}n^2
\right) +C +\Theta(n)}&
(\mbox{by Lemma 2})\\[.25in]
&=&{\displaystyle a n\log n -\frac{a}{ 4} n +  C +\Theta(n)}\\[.25in]
&=&{\displaystyle a n\log n +
\left( \Theta(n) + C -\frac{a}{ 4} n\right) }\\[.25in]
&\leq&{\displaystyle a n\log n}&(\mbox{choose } a \mbox{ so }
\Theta(n) + C \leq \frac{a}{ 4} n)\\
\end{array}
$$
We have shown that with an appropriate choice of $a$, $T(n)\leq a n \log n$
for all $n\geq 2$, so 
$T(n)=O(n\log n)$.

\noindent
We leave it to the reader to show that
$T(n)=\Omega(n\log n)$.  
\end{pf}
\end{thm}

\newpage
